        %%******************************************%%
        %%                                          %%
        %%        Modello di tesi di laurea         %%
        %%            di Andrea Giraldin            %%
        %%                                          %%
        %%             2 novembre 2012              %%
        %%                                          %%
        %%******************************************%%


% I seguenti commenti speciali impostano:
% 1. 
% 2. PDFLaTeX come motore di composizione;
% 3. tesi.tex come documento principale;
% 4. il controllo ortografico italiano per l'editor.

% !TEX encoding = UTF-8
% !TEX TS-program = pdflatex
% !TEX root = tesi.tex
% !TEX spellcheck = it-IT

\documentclass[10pt,                    % corpo del font principale
               a4paper,                 % carta A4
               twoside,                 % impagina per fronte-retro
               openright,               % inizio capitoli a destra
               english,                 
               italian,                 
               ]{book}    

\usepackage[utf8]{inputenc}             % codifica di input; anche [latin1] va bene
                                        % NOTA BENE! va accordata con le preferenze dell'editor

%**************************************************************
% Importazione package
%************************************************************** 

%\usepackage{amsmath,amssymb,amsthm}    % matematica

\usepackage[a4paper,                    % carta A4
            top=3cm,                    % margine superiore
            bottom=3cm,                 % margine inferiore
            head=30pt                   % altezza degli header
            ]{geometry}                 


\usepackage[english, italian]{babel}    % per scrivere in italiano e in inglese;
                                        % l'ultima lingua (l'italiano) risulta predefinita

\usepackage{bookmark}                   % segnalibri

\usepackage{caption}                    % didascalie

%\usepackage{chngpage,calc}              % centra il frontespizio

\usepackage{fancyhdr}                   % Gestisce gli header

\usepackage{csquotes}                   % gestisce automaticamente i caratteri (")

\usepackage{emptypage}                  % pagine vuote senza testatina e piede di pagina

\usepackage{epigraph}					% per epigrafi

\usepackage{eurosym}                    % simbolo dell'euro

\usepackage[T1]{fontenc}                % codifica dei font:
                                        % NOTA BENE! richiede una distribuzione *completa* di LaTeX

%\usepackage{indentfirst}               % rientra il primo paragrafo di ogni sezione

\usepackage{graphicx}                   % immagini

\usepackage{hyperref}                   % collegamenti ipertestuali



\usepackage[binding=5mm]{layaureo}      % margini ottimizzati per l'A4; rilegatura di 5 mm

\usepackage{listings}                   % codice

\usepackage{microtype}                  % microtipografia

\usepackage{mparhack,fixltx2e,relsize}  % finezze tipografiche

\usepackage{nameref}                    % visualizza nome dei riferimenti                                      

\usepackage[font=small]{quoting}        % citazioni

\usepackage{subfig}                     % sottofigure, sottotabelle

\usepackage[italian]{varioref}          % riferimenti completi della pagina

\usepackage[dvipsnames]{xcolor}         % colori

\usepackage{booktabs}                   % tabelle                                       
\usepackage{tabularx}                   % tabelle di larghezza prefissata                                    
\usepackage{longtable}                  % tabelle su più pagine                                        
\usepackage{ltxtable}                   % tabelle su più pagine e adattabili in larghezza

\usepackage[toc, acronym]{glossaries}   % glossario
                                        % per includerlo nel documento bisogna:
                                        % 1. compilare una prima volta tesi.tex;
                                        % 2. eseguire: makeindex -s tesi.ist -t tesi.glg -o tesi.gls tesi.glo
                                        % 3. eseguire: makeindex -s tesi.ist -t tesi.alg -o tesi.acr tesi.acn
                                        % 4. compilare due volte tesi.tex.

\usepackage[backend=biber,style=verbose-ibid,hyperref,backref]{biblatex}
                                        % eccellente pacchetto per la bibliografia; 
                                        % produce uno stile di citazione autore-anno; 
                                        % lo stile "numeric-comp" produce riferimenti numerici
                                        % per includerlo nel documento bisogna:
                                        % 1. compilare una prima volta tesi.tex;
                                        % 2. eseguire: biber tesi
                                        % 3. compilare ancora tesi.tex.

%**************************************************************
% file contenente le impostazioni della tesi
%**************************************************************

%**************************************************************
% Frontespizio
%**************************************************************

% Autore
\newcommand{\myName}{Federico Silvio Busetto}       
\newcommand{\matricola}{1026925}               

% Azienda
\newcommand{\myCompany}{IKS s.r.l.}
% Titolo della tesi              
\newcommand{\myTitle}{Un sistema di continuous building basato su Jenkins e tecnologie di containerizzazione}

% Tipo di tesi                   
\newcommand{\myDegree}{Tesi di laurea triennale}

% Università             
\newcommand{\myUni}{Università degli Studi di Padova}

% Facoltà       
\newcommand{\myFaculty}{Corso di Laurea in Informatica}

% Dipartimento
\newcommand{\myDepartment}{Dipartimento di Matematica "Tullio Levi-Civita"}

% Titolo del relatore
\newcommand{\profTitle}{Prof.}

% Relatore  e tutor aziendale
\newcommand{\myProf}{Gilberto Filè}
\newcommand{\myTutor}{Massimo Celegato}

% Luogo
\newcommand{\myLocation}{Padova}

% Anno accademico
\newcommand{\myAA}{2016-2017}

% Data discussione
\newcommand{\myTime}{07 Dicembre 2017}


%**************************************************************
% Impostazioni di impaginazione
% see: http://wwwcdf.pd.infn.it/AppuntiLinux/a2547.htm
%**************************************************************

\pagestyle{fancy}
\fancyhf{}
\setlength{\headheight}{1cm} 
\setlength{\parindent}{0pt}   % larghezza rientro della prima riga
\setlength{\parskip}{0pt}     % distanza tra i paragrafi
\renewcommand{\chaptermark}[1]{\markboth{#1}{}} 


\fancypagestyle{IHA-fancy-style}{%
    \fancyhf{}% Clear header and footer
    \fancyhead[LE,RO]{\slshape \rightmark}
    \fancyhead[LO,RE]{\slshape \leftmark}
    \fancyfoot[LE,RO]{\thepage} % Left side on Even pages; Right side on Odd pages
    \renewcommand{\headrulewidth}{0.4pt}% Line at the header visible
    \renewcommand{\footrulewidth}{0.4pt}% Line at the footer visible
}

% Redefine the plain page style
\fancypagestyle{plain}{%
   \fancyhf{}%
   \fancyfoot[LE,RO]{\thepage} % Left side on Even pages; Right side on Odd pages
   \renewcommand{\headrulewidth}{0pt}% Line at the header invisible
   \renewcommand{\footrulewidth}{0.4pt}% Line at the footer visible
}

%***INSERIMENTO DI NUOVE SOTTOSEZIONI
\setcounter{secnumdepth}{5}		% mostra nel documento fino al livello 8 (1.2.3.4)
\setcounter{tocdepth}{5}			% mostra nell'indice fino al livello 8 (1.2.3.4)

%**************************************************************
% Impostazioni di biblatex
%**************************************************************
\bibliography{bibliografia} % database di biblatex 

\defbibheading{bibliography} {
    \cleardoublepage
    \phantomsection 
    \addcontentsline{toc}{chapter}{\bibname}
    \chapter*{\bibname\markboth{\bibname}{\bibname}}
}

\setlength\bibitemsep{1.5\itemsep} % spazio tra entry

\DeclareBibliographyCategory{opere}
\DeclareBibliographyCategory{web}

\addtocategory{opere}{womak:lean-thinking}
\addtocategory{web}{site:agile-manifesto}

\defbibheading{opere}{\section*{Riferimenti bibliografici}}
\defbibheading{web}{\section*{Siti Web consultati}}


%**************************************************************
% Impostazioni di caption
%**************************************************************
\captionsetup{
    tableposition=top,
    figureposition=bottom,
    font=small,
    format=hang,
    labelfont=bf
}

%**************************************************************
% Impostazioni di glossaries
%**************************************************************

%**************************************************************
% Acronimi
%**************************************************************
\renewcommand{\acronymname}{Acronimi e abbreviazioni}

\newacronym[description={\glslink{apig}{Application Program Interface}}]
    {api}{API}{Application Program Interface}

\newacronym[description={\glslink{bash}{Burn Again Shell}}]
    {bash}{Bash}{Burn Again Shell}

\newacronym[description={\glslink{CentOS}{Community enterprise Operating System}}]
    {CentOS}{CentOS}{Community enterprise Operating System}
    
\newacronym[description={\glslink{CIg}{Configuration Item}}]
    {CI}{CI}{Configuration Item}

\newacronym[description={\glslink{gnug}{GNU is Not Unix}}]
    {GNU}{GNU}{}

\newacronym[description={\glslink{ide}{Integrated Development Environment}}]
{ide}{IDE}{Integrated Development Environment}

\newacronym[description={\glslink{iot}{Internet of Things}}]
{iot}{IoT}{Internet of Things}

\newacronym[description={\glslink{jdk}{Java Development Kit}}]
    {jdk}{JDK}{Java Development Kit}
    
\newacronym[description={\glslink{json}{Javascript Object Notation}}]
    {json}{JSON}{Javascript Object Notation}
    
\newacronym[description={\glslink{jvm}{Java Virtual Machine}}]
{jvm}{JVM}{Java Virtual Machine}

\newacronym[description={\glslink{pom}{Project Object Model}}]
    {pom}{POM}{Project Object Model}

\newacronym[description={\glslink{umlg}{Unified Modeling Language}}]
    {uml}{UML}{Unified Modeling Language}
    
\newacronym[description={\glslink{xml}{eXtensible Markup Language}}]
    {xml}{XML}{eXtensible Markup Language}

%**************************************************************
% Glossario
%**************************************************************
%\renewcommand{\glossaryname}{Glossario}

\newglossaryentry{apig}
{
    name=\glslink{api}{API},
    text=Application Program Interface,
    sort=api,
    description={in informatica con il termine \emph{Application Programming Interface API} (ing. interfaccia di programmazione di un'applicazione) si indica ogni insieme di procedure disponibili al programmatore, di solito raggruppate a formare un set di strumenti specifici per l'espletamento di un determinato compito all'interno di un certo programma. La finalità è ottenere un'astrazione, di solito tra l'hardware e il programmatore o tra software a basso e quello ad alto livello semplificando così il lavoro di programmazione}
}

\newglossaryentry{umlg}
{
    name=\glslink{uml}{UML},
    text=UML,
    sort=uml,
    description={In ingegneria del software \emph{UML, Unified Modeling Language} (ing. linguaggio di modellazione unificato) è un linguaggio di modellazione e specifica basato sul paradigma object-oriented. L'\emph{UML} svolge un'importantissima funzione di ``lingua franca'' nella comunità della progettazione e programmazione a oggetti. Gran parte della letteratura di settore usa tale linguaggio per descrivere soluzioni analitiche e progettuali in modo sintetico e comprensibile a un vasto pubblico}
}

\newglossaryentry{agent}
{
    name=\glslink{agent}{agent},
    text=agent,
    sort=agent,
    description={Sinonimo di slave. Nodo fisico, virtuale o container, che esegue compiti definiti e diretti da un'istanza \textit{master}}
}

\newglossaryentry{agile}
{
    name=\glslink{agile}{agile},
    text=agile,
    sort=agile,
    description={Insieme di metodi di sviluppo del software emersi a partire dai primi anni 2000, fondati su un insieme di principi comuni, direttamente o indirettamente derivati dai princìpi del \textit{Manifesto per lo sviluppo agile del software  \footnote{\url{http://agilemanifesto.org/}}}}
}

\newglossaryentry{application-server}
{
    name=\glslink{application-server}{application-server},
    text=Application Server,
    sort=application server,
    description={Tipologia di server che fornisce l'infrastruttura e le funzionalità di supporto, sviluppo ed esecuzione di applicazioni nonché altri componenti server in un contesto distribuito}
}

\newglossaryentry{artefatto}
{
    name=\glslink{artefatto}{artefatto},
    text=artefatto,
    sort=artefatto,
    description={Un prodotto immutabile, generato durante una \textit{build} che verrà opportunamente archiviato in un \textit{repository} e reso disponibile a determinati utenti per usi specifici (es. rilascio, testing)}
}

\newglossaryentry{baseline}
{
    name=\glslink{baseline}{baseline},
    text=baseline,
    sort=baseline,
    description={Nel ciclo di vita di un progetto, punto d'arrivo tecnico dal quale non si retrocede, sul quale calcolare l'avanzamento del lavoro in un progetto}
}

\newglossaryentry{blockchain}
{
    name=\glslink{blockchain}{blockchain},
    text=blockchain,
    sort=blockchain,
    description={Base di dati distribuita, introdotta dalla valuta Bitcoin che mantiene in modo continuo una lista crescente di record, i quali fanno riferimento a record precedenti presenti nella lista stessa in modo da resistere a manomissioni}
}

\newglossaryentry{build}
{
    name=\glslink{build}{build},
    text=build,
    sort=build,
    description={La costruzione, osservabile e tangibile di una componente software a partire dal suo codice sorgente.\\ 
    Risultato della singola esecuzione di un \textit{project}, in \textit{Jenkins}}
}

\newglossaryentry{bytecode}
{
    name=\glslink{bytecode}{bytecode},
    text=bytecode,
    sort=bytecode,
    description={Linguaggio intermedio tra il codice sorgente e il linguaggio macchina}
}

\newglossaryentry{CIg}
{
    name=\glslink{CI}{CI},
    text=Configuration Item,
    sort=configuration item,
    description={Singolo elemento sottoposto ad attività di gestione di configurazione. Ogni CI è identificato univocamente da:
        \begin{itemize}
            \item Codice Identificativo (ID);
            \item Nome;
            \item Data;
            \item Autore;
            \item Registro delle modifiche;
            \item Stato corrente;
        \end{itemize}}
}

\newglossaryentry{cloud}
{
    name=\glslink{cloud}{cloud},
    text=cloud,
    sort=cloud,
    description={Abbreviazione di cloud computing. Paradigma informatico per l'utilizzo di risorse hardware o software in remoto, tramite la disponibilità on demand di risorse preesistenti e configurabili.}
}

\newglossaryentry{container}
{
    name=\glslink{container}{container},
    text=container,
    sort=container,
    description={Processo isolato dal resto del sistema che rappresenta l'istanza in esecuzione di una data immagine, secondo opportune configurazioni aggiuntive}
}

\newglossaryentry{continuous building}
{
    name=\glslink{continuous building}{continuous building},
    text=continuous building,
    sort=continuous building,
    description={In ingegneria del software per \emph{continuous building}, (spesso identificata come ``Build automation'') si intende il processo di automazione dell'intero ciclo di build a partire dalla compilazione del codice sorgente, al garantire che i risultati siano sempre ripetibili e conformi alle attese, all'esecuzione dei test e al rendere disponibile gli artefatti risultanti. Tale processo è una parte consistente dell'\textit{integrazione continua}}
}

\newglossaryentry{debugger}
{
    name=\glslink{debugger}{debugger},
    text=debugger,
    sort=debugger,
    description={Strumento per agevolare l'individuazione e permettere un'agevole rimozione di errori software, comunemente chiamati bug}
}

\newglossaryentry{deploy}
{
    name=\glslink{deploy}{deploy},
    text=deploy,
    sort=deploy,
    description={L'insieme delle attività che permettono ad una componente software di essere disponibile ad un determinato uso}
}

\newglossaryentry{gnug}
{
    name=\glslink{GNU}{GNU},
    text=GNU,
    sort=gnu,
    description={Sistema operativo Unix-like, ideato nel 1984 da Richard Stallman e promosso dalla Free Software Foundation, allo scopo di ottenere un sistema operativo completo utilizzando esclusivamente software libero}
}

\newglossaryentry{integrazione-continua}
{
    name=\glslink{integrazione-continua}{integrazione-continua},
    text=integrazione continua,
    sort=integrazione continua,
    description={Insieme delle attività che permettono un allineamento frequente degli ambienti di lavoro di più sviluppatori}
}

\newglossaryentry{immagine}
{
    name=\glslink{immagine}{immagine},
    text=immagine,
    sort=immagine,
    description={Base fondante e immutabile utilizzata per la costruzione di \textit{container}}
}
 
\newglossaryentry{Java EE}
{
    name=\glslink{Java EE}{Java EE},
    text=Java Enterprise Edition,
    sort=Java EE,
    description={Insieme di specifiche le cui implementazioni vengono sviluppate in linguaggio di programmazione Java con ampio utilizzo nella programmazione Web}
}

\newglossaryentry{Linux}
{
    name=\glslink{Linux}{Linux},
    text=Linux,
    sort=linux,
    description={Famiglia di sistemi operativi di tipo Unix-like, pubblicati sotto varie possibili distribuzioni, aventi la caratteristica comune di utilizzare come nucleo il kernel Linux}
}

\newglossaryentry{nodo}
{
    name=\glslink{nodo}{nodo},
    text=nodo,
    sort=nodo,
    description={Insieme di specifiche le cui implementazioni vengono sviluppate in linguaggio di programmazione Java con ampio utilizzo nella programmazione Web}
}

\newglossaryentry{machine-learning}
{
    name=\glslink{machine-learning}{machine-learning},
    text=machine-learning,
    sort=machine-learning,
    description={Campo dell'informatica che fornisce alle macchine l'abilità di imparare, nello svolgere determinati compiti, senza essere programmate esplicitamente}
}

\newglossaryentry{master}
{
    name=\glslink{master}{master},
    text=master,
    sort=master,
    description={L'unità centrale di coordinazione dei processi che possiede la configurazione, carica i plugin e gestisce tutti gli aspetti di esecuzione dei vari \textit{projects}}
}

\newglossaryentry{middleware}
{
    name=\glslink{middleware}{middleware},
    text=middleware,
    sort=middleware,
    description={Insieme di programmi informatici che fungono da intermediari tra diverse applicazioni e componenti software}
}

\newglossaryentry{plugin}
{
    name=\glslink{plugin}{plugin},
    text=plugin,
    sort=plugin,
    description={Componente software che fornisce funzionalità aggiuntive ad estensione di un determinato prodotto software preesistente ma separatamente da esso}
}

\newglossaryentry{project}
{
    name=\glslink{project}{project},
    text=project,
    sort=project,
    description={Una configurazione definita da un utente, che descrive un lavoro che Jenkins dovrà svolgere, come ad esempio costruire una componente software, eseguire uno script, ecc... A volte riferito come \textit{job}}
}

\newglossaryentry{Red Hat}
{
    name=\glslink{Red Hat}{Red Hat},
    text=Red Hat,
    sort=red hat,
    description={Società multinazionale statunitense che si dedica allo sviluppo software e al supporto di software libero e open source in ambiente aziendale}
}

\newglossaryentry{repository}
{
    name=\glslink{repository}{repository},
    text=repository,
    sort=repository,
    description={Base di dati centralizzata nella quale risiedono, individualmente, tutti i \gls{CI} di ogni \gls{baseline} nella loro storia completa}
}

\newglossaryentry{slave}
{
    name=\glslink{slave}{slave},
    text=slave,
    sort=slave,
    description={Letteralmente, schiavo. In Jenkins rappresenta un nodo sul quale vengono eseguiti più compiti, definiti sotto forma di \textit{project}, secondo la direzione di un \textit{master}. A volte definito come \textit{Agent}}
}

\newglossaryentry{system integration}
{
    name=\glslink{system integration}{system integration},
    text=system integration,
    sort=system integration,
    description={Processo che consiste nel collegare assieme differenti sistemi e applicazioni software, a livello fisico o funzionale in modo che agiscano come un insieme coordinato e coeso}
} % database di termini
\makeglossaries


%**************************************************************
% Impostazioni di graphicx
%**************************************************************
\graphicspath{{immagini/}} % cartella dove sono riposte le immagini


%**************************************************************
% Impostazioni di hyperref
%**************************************************************
\hypersetup{
    %hyperfootnotes=false,
    %pdfpagelabels,
    %draft,	% = elimina tutti i link (utile per stampe in bianco e nero)
    colorlinks=true,
    linktocpage=true,
    pdfstartpage=1,
    pdfstartview=FitV,
    % decommenta la riga seguente per avere link in nero (per esempio per la stampa in bianco e nero)
    %colorlinks=false, linktocpage=false, pdfborder={0 0 0}, pdfstartpage=1, pdfstartview=FitV,
    breaklinks=true,
    pdfpagemode=UseNone,
    pageanchor=true,
    pdfpagemode=UseOutlines,
    plainpages=false,
    bookmarksnumbered,
    bookmarksopen=true,
    bookmarksopenlevel=1,
    hypertexnames=true,
    pdfhighlight=/O,
    %nesting=true,
    %frenchlinks,
    urlcolor=webbrown,
    linkcolor=RoyalBlue,
    citecolor=webgreen,
    %pagecolor=RoyalBlue,
    %urlcolor=Black, linkcolor=Black, citecolor=Black, %pagecolor=Black,
    pdftitle={\myTitle},
    pdfauthor={\textcopyright\ \myName, \myUni, \myFaculty},
    pdfsubject={},
    pdfkeywords={},
    pdfcreator={pdfLaTeX},
    pdfproducer={LaTeX}
}

%**************************************************************
% Impostazioni di itemize
%**************************************************************
\renewcommand{\labelitemi}{$\ast$}

%\renewcommand{\labelitemi}{$\bullet$}
%\renewcommand{\labelitemii}{$\cdot$}
%\renewcommand{\labelitemiii}{$\diamond$}
%\renewcommand{\labelitemiv}{$\ast$}


%**************************************************************
% Impostazioni di listings
%**************************************************************

%Dockerfile
\lstdefinelanguage{docker}{
    keywords={FROM, RUN, COPY, ADD, ENTRYPOINT, CMD,  ENV, WORKDIR, EXPOSE, LABEL, USER, VOLUME, STOPSIGNAL, ONBUILD, MAINTAINER},
    keywordstyle=\color{blue}\bfseries,
    identifierstyle=\color{black},
    sensitive=false,
    comment=[l]{\#},
    numbers=left,
    stepnumber=1,
    commentstyle=\color{purple}\ttfamily,
    stringstyle=\color{red}\ttfamily,
    morestring=[b]',
    morestring=[b]"
}

%C++
\lstset{
    language=[LaTeX]Tex,
    keywordstyle=\color{RoyalBlue}, %\bfseries,
    basicstyle=\small\ttfamily,
    %identifierstyle=\color{NavyBlue},
    commentstyle=\color{Green}\ttfamily,
    stringstyle=\rmfamily,
    numbers=none, %left,%
    numberstyle=\scriptsize, %\tiny
    stepnumber=5,
    numbersep=8pt,
    showstringspaces=false,
    breaklines=true,
    frameround=ftff,
    frame=single
} 


%**************************************************************
% Impostazioni di xcolor
%**************************************************************
\definecolor{webgreen}{rgb}{0,.5,0}
\definecolor{webbrown}{rgb}{.6,0,0}
\definecolor{mygreen}{rgb}{0,0.6,0}
\definecolor{darkgreen}{RGB}{0,100,0}
\definecolor{mygray}{rgb}{0.5,0.5,0.5}
\definecolor{mymauve}{rgb}{0.58,0,0.82}
\definecolor{lightGray}{RGB}{150, 150, 150}
\definecolor{darkblue}{rgb}{0.07, 0.04, 0.56}
\definecolor{lightblue}{RGB}{176, 196, 222}
\definecolor{tableHeaderRed}{RGB}{211, 47, 0}
\definecolor{tableHeaderBlue}{RGB}{2, 49, 103}
\definecolor{tableLineOne}{RGB}{245, 245, 245}
\definecolor{tableLineTwo}{RGB}{224, 224, 224}
\definecolor{Maroon}{rgb}{0.48, 0.07, 0.07}

%**************************************************************
% Impostazioni quote
%**************************************************************

\newcommand*\openquote{\makebox(25,-22){\scalebox{5}{``}}}
\newcommand*\closequote{\makebox(25,-22){\scalebox{5}{''}}}
\colorlet{shadecolor}{Azure}

\makeatletter
\newif\if@right
\def\shadequote{\@righttrue\shadequote@i}
\def\shadequote@i{\begin{snugshade}\begin{quote}\openquote}
        \def\endshadequote{%
            \if@right\hfill\fi\closequote\end{quote}\end{snugshade}}
\@namedef{shadequote*}{\@rightfalse\shadequote@i}
\@namedef{endshadequote*}{\endshadequote}
\makeatother

%**************************************************************
% Altro
%**************************************************************

\newcommand{\omissis}{[\dots\negthinspace]} % produce [...]

% eccezioni all'algoritmo di sillabazione
\hyphenation
{
    ma-cro-istru-zio-ne
}

\newcommand{\sectionname}{sezione}
\addto\captionsitalian{\renewcommand{\figurename}{Figura}
                       \renewcommand{\tablename}{Tabella}}

\newcommand{\glsfirstoccur}{\ap{{[g]}}}

\newcommand{\intro}[1]{\emph{\textsf{#1}}}

%**************************************************************
% Environment per ``rischi''
%**************************************************************
\newcounter{riskcounter}                % define a counter
\setcounter{riskcounter}{0}             % set the counter to some initial value

%%%% Parameters
% #1: Title
\newenvironment{risk}[1]{
    \refstepcounter{riskcounter}        % increment counter
    \par \noindent                      % start new paragraph
    \textbf{\arabic{riskcounter}. #1}   % display the title before the 
                                        % content of the environment is displayed 
}{
    \par\medskip
}

\newcommand{\riskname}{Rischio}

\newcommand{\riskdescription}[1]{\textbf{\\Descrizione:} #1.}

\newcommand{\risksolution}[1]{\textbf{\\Soluzione:} #1.}

%**************************************************************
% Environment per ``use case''
%**************************************************************
\newcounter{usecasecounter}             % define a counter
\setcounter{usecasecounter}{0}          % set the counter to some initial value

%%%% Parameters
% #1: ID
% #2: Nome
\newenvironment{usecase}[2]{
    \renewcommand{\theusecasecounter}{\usecasename #1}  % this is where the display of 
                                                        % the counter is overwritten/modified
    \refstepcounter{usecasecounter}             % increment counter
    \vspace{10pt}
    \par \noindent                              % start new paragraph
    {\large \textbf{\usecasename #1: #2}}       % display the title before the 
                                                % content of the environment is displayed 
    \medskip
}{
    \medskip
}

\newcommand{\usecasename}{UC}

\newcommand{\usecaseactors}[1]{\textbf{\\Attori Principali:} #1. \vspace{4pt}}
\newcommand{\usecasepre}[1]{\textbf{\\Precondizioni:} #1. \vspace{4pt}}
\newcommand{\usecasedesc}[1]{\textbf{\\Descrizione:} #1. \vspace{4pt}}
\newcommand{\usecasepost}[1]{\textbf{\\Postcondizioni:} #1. \vspace{4pt}}
\newcommand{\usecasealt}[1]{\textbf{\\Scenario Alternativo:} #1. \vspace{4pt}}

%**************************************************************
% Environment per ``namespace description''
%**************************************************************

\newenvironment{namespacedesc}{
    \vspace{10pt}
    \par \noindent                              % start new paragraph
    \begin{description} 
}{
    \end{description}
    \medskip
}

\newcommand{\classdesc}[2]{\item[\textbf{#1:}] #2}                     % file con le impostazioni personali

\begin{document}
%**************************************************************
% Materiale iniziale
%**************************************************************
\frontmatter
\begin{titlepage}
    
    % Defines a new command for the horizontal lines, change thickness here
    \newcommand{\HRule}{\rule{\linewidth}{0.5mm}} 
    
    % Center everything
    \center
    
    %----------------------------------------------------------------------------------------
    %	HEADING SECTIONS
    %----------------------------------------------------------------------------------------
    
    % Name of your university/college
    
    \begin{LARGE}
        \textsc{\textbf{\myUni}}\\[1cm] 
    \end{LARGE}
    
    %----------------------------------------------------------------------------------------
    %	LOGO SECTION
    %----------------------------------------------------------------------------------------
    
    \includegraphics[height=6cm]{immagini/logo-unipd.png}\\[1cm]
    
    \includegraphics[height=1.5cm, width = 9.25cm]{immagini/MathDip.png}\\
    \textsc{\myDepartment}\\[1cm]
    
    % Major heading such as course name
    \textsc{\Large Scuola di Scienze}\\[0.5cm] 
    
    % Minor heading such as course title
    \textsc{\large Corso di Laurea in Informatica}\\[0.5cm] 
    
    \vspace{1cm}
    
    %----------------------------------------------------------------------------------------
    %	TITLE SECTION
    %----------------------------------------------------------------------------------------
    
    \begin{LARGE}
        \begin{center}
            \textbf{\myTitle}\\
        \end{center}
    \end{LARGE}
    
    %----------------------------------------------------------------------------------------
    %	AUTHOR SECTION
    %----------------------------------------------------------------------------------------
    
    \vspace{1.5cm}
    
    \begin{large}
        \begin{flushleft}
            \textit{Relatore}\\ 
            \vspace{5pt} 
            \profTitle \myProf
        \end{flushleft}
        
        \begin{flushright}
            \textit{Laureando}\\ 
            \vspace{5pt} 
            \myName\\
            \matricola
        \end{flushright}
    \end{large}
    \vspace{1.3cm}
    
    %----------------------------------------------------------------------------------------
    %	DATE SECTION
    %----------------------------------------------------------------------------------------
    \line(1, 0){385} \\
    \begin{normalsize}
        \textsc{Anno Accademico \myAA}
    \end{normalsize}
    
    %----------------------------------------------------------------------------------------
    % Fill the rest of the page with whitespace
    \vfill 
    
\end{titlepage}
% !TEX encoding = UTF-8
% !TEX TS-program = pdflatex
% !TEX root = ../tesi.tex

%**************************************************************
% Colophon
%**************************************************************
\clearpage
\phantomsection
\thispagestyle{empty}

\hfill

\vfill

\noindent\myName: \textit{\myTitle,}
\myDegree,
\textcopyright\ \myTime.
% !TEX encoding = UTF-8
% !TEX TS-program = pdflatex
% !TEX root = ../tesi.tex

%**************************************************************
% Dedica
%**************************************************************
\cleardoublepage
\phantomsection
\thispagestyle{empty}
\pdfbookmark{Dedica}{Dedica}

\vspace*{3cm}

\begin{center}
Solo chi non sa dove sta andando può arrivare lontano. \\ \medskip
--- Oscar Wilde    
\end{center}

\medskip

\begin{center}
Dedicato a mio nonno Gigi.
\end{center}

% !TEX encoding = UTF-8
% !TEX TS-program = pdflatex
% !TEX root = ../tesi.tex

%**************************************************************
% Sommario
%**************************************************************
\cleardoublepage
\phantomsection
\pdfbookmark{Sommario}{Sommario}
\begingroup
\let\clearpage\relax
\let\cleardoublepage\relax
\let\cleardoublepage\relax

\chapter*{Sommario}

Il presente documento descrive il lavoro svolto durante il periodo di stage, della durata di circa trecento ore, dal laureando \myName\ presso l'azienda \myCompany, sita in Padova.\\

Lo scopo principale dello stage consisteva nella realizzazione di un sistema di \textit{\gls{continuous building}}, in grado di prendere in ingresso il codice sorgente di un'applicazione d'esempio, effettuarne la compilazione, e produrre in uscita un \gls{artefatto} pronto per il \textit{\gls{deploy}}.

Era inoltre richiesto che tale sistema prevedesse l'utilizzo di tecnologie di containerizzazione, al fine di valutare gli effettivi vantaggi della loro adozione.

Particolare attenzione è stata posta agli aspetti legati all'esercibilità, quali: monitoraggio, aggiornamento, backup e ripristino della piattaforma.\\

Gli obiettivi da raggiungere erano molteplici:\\
In primo luogo era richiesta la progettazione e la definizione dell'architettura sia di alto livello che di dettaglio di tale sistema, e la sua implementazione.

In secondo luogo è stato richiesto uno studio di fattibilità sulle modalità di comunicazione tra nodi, approfondendo in particolare le comunicazioni tra nodi container in un ambiente distribuito.

Infine era richiesto lo sviluppo di un \textit{\gls{plugin}} base per Jenkins, da utilizzare come modello di riferimento per l'eventuale implementazione di nuove funzionalità a livello di piattaforma.

%\vfill
%
%\selectlanguage{english}
%\pdfbookmark{Abstract}{Abstract}
%\chapter*{Abstract}
%
%\selectlanguage{italian}

\endgroup			

\vfill


% !TEX encoding = UTF-8
% !TEX TS-program = pdflatex
% !TEX root = ../tesi.tex

%**************************************************************
% Ringraziamenti
%**************************************************************
\cleardoublepage
\phantomsection
\pdfbookmark{Ringraziamenti}{ringraziamenti}

\begin{flushright}{
	\slshape    
	``Una grande illuminazione nasce da un grande dubbio.''} \\ 
	\medskip
    --- Proverbio Zen
\end{flushright}


\bigskip

\begingroup
\let\clearpage\relax
\let\cleardoublepage\relax
\let\cleardoublepage\relax

\chapter*{Ringraziamenti}

\noindent \textit{Voglio esprimere la mia gratitudine al Prof. \myProf, relatore della mia tesi, per la disponibilità, l'aiuto e il sostegno fornitomi durante il periodo di stage e nella stesura di questo documento. Un grazie sincero anche per l'attenzione che riserva ai suoi studenti, fornendo le video lezioni del corso di Programmazione.}\\
\bigskip

\noindent \textit{Un particolare ringraziamento al Prof. Tullio Vardanega, per aver dato una svolta al mio percorso accademico, insegnandomi a fissare i miei obiettivi e a pianificare adeguate strategie nel perseguirli.}\\
\bigskip

\noindent \textit{Desidero ringraziare con affetto i miei genitori e la mia famiglia per tutto l'amore che mi hanno donato e per aver sempre creduto in me, anche quando io stesso ne dubitavo.}\\
\bigskip

\noindent \textit{Un ringraziamento speciale va ai miei nonni e a mia zia Dirce, per avermi insegnato a non mollare mai e ad andare sempre avanti, anche quando sembra tutto sia perduto.}\\
\bigskip

\noindent \textit{Ho desiderio di ringraziare poi i miei compagni universitari, in particolar modo i ragazzi (e le ragazze) del Visions Team per aver reso quest'ultimo anno il migliore. Un grazie doppio ai miei amici Ettore, Thomas, Elisa, Marco, Paolo e Beatrice.}\\
\bigskip

\noindent \textit{Infine, ma non meno importanti, un ringraziamento ai colleghi di \myCompany, specialmente al mio tutor aziendale \myTutor\ per la pazienza e la disponibilità nei miei confronti e per avermi fatto crescere professionalmente; grazie ad Andrei per avermi accompagnato alla scoperta del fantastico mondo dei \textit{container} e un grazie sentito a Stefano, Daniele e Giulia per tutti i bei momenti trascorsi assieme.}\\
\bigskip

\noindent\textit{\myLocation, \myTime}
\hfill \myName

\endgroup


% !TEX encoding = UTF-8
% !TEX TS-program = pdflatex
% !TEX root = ../tesi.tex

%**************************************************************
% Indici
%**************************************************************
\cleardoublepage
\pdfbookmark{\contentsname}{tableofcontents}
\setcounter{tocdepth}{2}
\tableofcontents
%\markboth{\contentsname}{\contentsname} 
\clearpage

\begingroup 
    \let\clearpage\relax
    \let\cleardoublepage\relax
    \let\cleardoublepage\relax
    %*******************************************************
    % Elenco delle figure
    %*******************************************************    
    \phantomsection
    \pdfbookmark{\listfigurename}{lof}
    \listoffigures

    \vspace*{8ex}

    %*******************************************************
    % Elenco delle tabelle
    %*******************************************************
    \phantomsection
    \pdfbookmark{\listtablename}{lot}
    \listoftables
        
    \vspace*{8ex}
\endgroup

\cleardoublepage

\cleardoublepage

%**************************************************************
% Materiale principale
%**************************************************************
\mainmatter
% !TEX encoding = UTF-8
% !TEX TS-program = pdflatex
% !TEX root = ../tesi.tex

%**************************************************************
\pagestyle{IHA-fancy-style}
\chapter{Introduzione}
\label{cap:introduzione}
%**************************************************************

In questo capitolo introduttivo viene fornita una panoramica dell’azienda \myCompany, proponente dello stage da me svolto e l’idea alla base del progetto.
Vengono inoltre fissate alcune convenzioni tipografiche per agevolare la lettura del documento e viene presentata la struttura dello stesso, secondo una suddivisione per capitoli.

%\noindent Esempio di utilizzo di un termine nel glossario \\
%\gls{api}. \\

%\noindent Esempio di citazione in linea \\
%\cite{site:agile-manifesto}. \\

%\noindent Esempio di citazione nel pie' di pagina \\
%citazione\footcite{womak:lean-thinking} \\

%**************************************************************
\section{L'azienda}
%**************************************************************
IKS \footnote{\url{https://www.iks.it/}} s.r.l. (Information Knowledge Supply) nasce nel 1999 a Padova, come azienda di consulenza e \textit{\gls{system integration}} con servizi rivolti ad una variegata tipologia di clienti quali banche e centri servizi, industria e pubblica amministrazione.\\

L'azienda punta da sempre su temi strategici in capo informatico, quali la sicurezza informatica, lo sviluppo di soluzioni software \textit{\gls{Java EE}}, l'automazione e l'innovazione. \\

Questo approccio portò IKS negli anni 2000 ad essere uno dei leader nel settore, nel nord-est italiano, con sedi a Padova, Milano e Trento e nel 2016 a formare, assieme ad altre tre aziende (Insirio, System Evolution e Kirey) il gruppo Kirey\footnote{\url{http://kireygroup.com/}}, con l'obiettivo di:

\begin{shadequote}
    Diventare il principale progetto italiano di integrazione e collaborazione fra imprese del mondo IT e assumere il ruolo di protagonista di riferimento di tutte le aziende che vogliono affrontare percorsi di innovazione digitale, riunendo eccellenze di settore, competenze verticali e partnership tecnologiche di alto livello.\par\emph{Sito internet Kirey Group}
\end{shadequote}


\vspace{2.5em}
\begin{figure}[H]
    \capstart
    \centering
    \includegraphics[height=2cm, keepaspectratio]{../immagini/IKS.png}
    \caption{Logo di IKS s.r.l.}
\end{figure}

L'azienda è strutturata in \textit{business units (BU)}, cioè divisioni specializzate in particolari compiti. Ad esempio: la BU \textit{\gls{middleware}} si occupa della gestione della configurazione di applicativi dei clienti, andando a fornire opportune impostazioni dell'infrastruttura secondo le richieste pervenute.

%**************************************************************
\section{L'idea}
\label{sec:idea}
%**************************************************************
Il mercato dello sviluppo software sta cambiando; mentre in passato si assisteva alla nascita di piccole medie imprese, focalizzate sullo sviluppo di uno o due prodotti, con tempi di rilascio lunghi tra una versione e l'altra, oggi la competizione è serrata. \\ Negli ultimi anni si assiste alla nascita di nuovi paradigmi: il \gls{machine-learning}, con macchine in grado di eseguire autonomamente alcune operazioni; l'\gls{iot}, con dispositivi perennemente connessi alla rete, in grado di misurare l'ambiente che li circonda e comunicare tra loro; e infine un'automazione sempre più spinta in tutti gli ambiti delle nostre vite quotidiane.\\
Da applicazioni monolitiche (un unico prodotto software dalle molte funzionalità, strettamente dipendenti e correlate tra loro), si sta passando allo sviluppo di soluzioni flessibili a microservizi, dove vari moduli con funzionalità univoche e distinte vengono composti con altri moduli, al fine di ottenere un'architettura scalabile, snella e robusta.\\

\begin{figure}[H]
    \capstart
    \captionsetup{justification=centering}
    \centering
    \includegraphics[width=9cm, keepaspectratio]{../immagini/Microservice_Architecture.png}
    \caption[Esempio d'applicazione a microservizi]{Esempio d'applicazione a microservizi, immagine tratta da \\ \url{http://microservices.io/patterns/microservices.html}}
\end{figure}

In quest'ottica, va ripensato l'intero ciclo di sviluppo del software; un'azienda, anche di grandi dimensioni, dovrà mettere in campo dei processi evolutivi per non rischiare di perdere posizioni di mercato. 

IKS fornisce ai suoi clienti degli appositi servizi di consulenza, accompagnandoli nella transizione verso una metodologia di sviluppo \textit{Devops}; tale metodologia di sviluppo è un'evoluzione del modello \gls{agile}, incentrato su un'ampia automazione delle operazioni ripetitive (quali ad esempio il testing e il rilascio) e sul principio cardine di collaborazione stretta tra sviluppo, ambito \textit{operations} (cioè gestione dell'infrastruttura, mantenimento dell'ambiente di lavoro, configurazione, ecc...) e garanzia della qualità, in passato visti come mondi separati e ben distinti. \\

Scopo dello stage è lo studio e la realizzazione di un sistema di \gls{continuous building}, che in ottica \textit{Devops}, permetta un'automazione dei processi di \gls{build}, \textit{test} e \textit{rilascio} del software con un'attenzione particolare agli aspetti di portabilità del sistema, anche in previsione di un funzionamento dello stesso in ambienti \gls{cloud}.

\begin{figure}[H]
    \capstart
    \centering
    \includegraphics[height=5cm, keepaspectratio]{../immagini/Devops.png}
    \caption[Visione d'insieme della metodologia \textit{Devops}]{Visione d'insieme della metodologia \textit{Devops}, immagine tratta da \\ \url{http://it.wikipedia.org}}
\end{figure}

%**************************************************************
\section{Organizzazione del documento}
%**************************************************************

Dopo questo primo capitolo introduttivo, il resto del documento si articola come segue:

\begin{description}
    
    \item[{\hyperref[cap:lo-stage]{Il secondo capitolo}}] descrive in dettaglio il progetto di stage, riportando le motivazioni che hanno spinto l'azienda a proporlo e le mie personali nella scelta, gli obiettivi da raggiungere e la relativa pianificazione.
    
    \item[{\hyperref[cap:tecnologie-adottate]{Il terzo capitolo}}] approfondisce le tecnologie adottate nella realizzazione del progetto, fornendo per ciascuna di esse una panoramica generale e contestualizzandone il ruolo nel sistema implementato.
    
    \item[{\hyperref[cap:realizzazione-del-sistema]{Il quarto capitolo}}] illustra le varie fasi di realizzazione del sistema, dalla progettazione all'effettiva implementazione.
    
    \item[{\hyperref[cap:vantaggi-e-punti-di-attenzione]{Il quinto capitolo}}]riporta una panoramica approfondita sui vantaggi e i punti di attenzione dell'architettura individuata.
    
    \item[{\hyperref[cap:esercibilità]{Il sesto capitolo}}]Illustra le procedure individuate in merito all'esercibilità del sistema, cioè agli aspetti di monitoraggio, upgrade, backup e ripristino della piattaforma.
        
    \item[{\hyperref[cap:verifica-validazione]{Il settimo capitolo}}] riporta le strategie di verifica e validazione adottate e il risultato dei test condotti.
    
    \item[{\hyperref[cap:conclusioni]{Nell'ottavo capitolo}}] infine, si riportano alcune considerazioni circa l'attività di stage e gli obiettivi raggiunti.
\end{description}

Riguardo la stesura del testo, relativamente al documento sono state adottate le seguenti convenzioni tipografiche:
\begin{itemize}
	\item gli acronimi, le abbreviazioni e i termini ambigui o di uso non comune menzionati vengono definiti nel glossario, situato alla fine del presente documento. \\
    Alcune definizioni dei termini del glossario, possono essere tratte da fonti online, tra cui: \\ 
    \url{https://it.wikipedia.org/};
	%\item Alla prima occorrenza dei termini riportati nel glossario viene utilizzata la seguente nomenclatura: \emph{parola}\glsfirstoccur;
	\item i termini in lingua straniera o facenti parti del gergo tecnico sono evidenziati con il carattere \emph{corsivo}. Fanno eccezione a questa norma, i nomi delle tecnologie.
\end{itemize}             % Introduzione
% !TEX encoding = UTF-8
% !TEX TS-program = pdflatex
% !TEX root = ../tesi.tex

%**************************************************************
\pagestyle{IHA-fancy-style}
\chapter{Lo stage}
\label{cap:lo-stage}
%**************************************************************

\intro{In questo capitolo viene illustrato dettagliatamente il progetto di stage, riportandone la pianificazione iniziale, gli obiettivi da raggiungere, i vincoli ai quali era sottoposto e le motivazioni personali nella scelta.
Viene inoltre riportata un'analisi preventiva dei possibili rischi e viene esposto l'ambiente di lavoro. }\\

%**************************************************************
\section{Introduzione al progetto}
%**************************************************************

Riprendendo il contesto definito nella sezione ``{\hyperref[sec:idea]{1.3 - L'idea}}'', occorre definire in modo più approfondito e dettagliato la metodologia di sviluppo \textit{Devops}; abbiamo visto come tale metodologia di sviluppo si focalizza sul rapido rilascio dei prodotti software e su un automazione spinta dei processi ripetibili. In tutto questo non vi è più una divisione netta e separata tra sviluppatori e l'ambito \textit{operations} ma viene instaurato un clima basato sulla collaborazione e sulla condivisione delle conoscenze. 

\begin{figure}[H]
    \capstart
    \centering
    \includegraphics[height=3cm, keepaspectratio]{../immagini/devopsbridge.png}
    \caption{Devops come punto d'unione tra sviluppo e operations, immagine tratta da \\ \url{https://intland.com/}}
\end{figure}

I processi di tale metodologia consistono nel:

\begin{itemize}
    \item Pianificare gli obiettivi da raggiungere e progettare le funzionalità da implementare nel prodotto software al fine che esso consegua tali obiettivi (\textit{Plan});
    \item Implementare effettivamente tali funzionalità (\textit{Code});
    \item Costruire le componenti software e i relativi artefatti (\textit{Build});
    \item Assicurarsi che le componenti software costruite siano conformi alle attese e ai bisogni (\textit{Test})
    \item Rilasciare il prodotto (\textit{Release});
    \item Rendere disponibile il prodotto a determinati usi (\textit{Deploy});
    \item Far si che il prodotto sia operativo, tramite opportune configurazioni (\textit{Operate});
    \item Monitorare le performance e lo stato del prodotto (\textit{Monitor}).
\end{itemize}

\begin{figure}[H]
    \capstart
    \centering
    \includegraphics[width=7cm, keepaspectratio]{../immagini/devops-process.png}
    \caption{I vari processi coinvolti all'interno della metodologia \textit{Devops}, immagine tratta da \\ \url{https://www.01net.it/devops-monitoring/}}
\end{figure}

Nel concreto, tali processi vengono svolti tramite una serie di attività e strumenti, tra cui:

\begin{itemize}
    \item Sviluppo del codice tramite opportuni ambienti di sviluppo integrato (\gls{ide}) e versionamento del codice che risiederà in un \gls{repository};
    \item Strumenti di \gls{integrazione-continua} al fine di automatizzare il processo di \textit{\gls{build}} e il loro monitoraggio;
    \item Attività continue e automatizzate di \textit{testing} al fine di verificare le funzionalità del prodotto, e garantire che esso abbia le desiderate proprietà qualitative; 
    \item Attività automatizzate di \textit{packaging}, con un \gls{repository} per gli artefatti;
    \item Attività di gestione dei cambiamenti, approvazione dei rilasci e automazione degli stessi;
    \item Strumenti di configurazione dell'infrastruttura, vista anch'essa come codice (\textit{Infrastructure as Code})
    \item Strumenti e metriche di monitoraggio continuo della \textit{user experience} e delle prestazioni del prodotto software.
\end{itemize}

Le sfide che tale metodologia si trova ad affrontare sono molteplici, ad esempio:

\begin{itemize}
    \item \textbf{Difficoltà nel mantenere e rendere disponibile l'ambiente di produzione}. A tal fine, si fa ampio uso di tecnologie di virtualizzazione e containerizzazione in modo da simulare gli ambienti, soprattutto durante le fasi di test dei prodotti;
    \item \textbf{Gestione dell'infrastruttura difficilmente automatizzabile}. Opportuni strumenti di gestione della configurazione permettono di gestire in modo proattivo l'infrastruttura, vista come codice, tramite linguaggi di configurazione facilmente personalizzabili;
    \item \textbf{Difficoltà nel diagnosticare mancanze e ricevere feedback sul prodotto.} Grazie a strumenti dedicati, è possibile un monitoraggio efficace e continuo delle applicazioni.  
\end{itemize}

Uno dei punti cardine è rappresentato dall'\textit{\gls{integrazione-continua}}, cioè l'allineamento continuo degli ambienti di lavoro di più sviluppatori. \\
Nello specifico l'\gls{integrazione-continua} si basa su:

\begin{itemize}
    \item Versionamento del codice;
    \item Automatizzazione delle build e dei test;
    \item Esecuzione dei test eseguiti in un clone dell'ambiente di produzione;
    \item Gestione e recupero degli artefatti;
    \item Automatizzazione dei rilasci.
\end{itemize} 

Il \textit{\gls{continuous building}}, facilmente confondibile con l'attività di \gls{integrazione-continua}
rappresenta una parte cospicua ma non il tutto. Non si occupa ad esempio della configurazione dei \gls{repository} per il versionamento del codice sorgente, ma assume che siano già presenti. Così come non si occupa direttamente della gestione dell'infrastruttura \gls{cloud}.

\begin{figure}[H]
    \capstart
    \centering
    \includegraphics[width=12cm, keepaspectratio]{../immagini/cb.png}
    \caption{In rosso, la parte di \gls{continuous building}, all'interno dell'\gls{integrazione-continua}, immagine tratta da \\ \url{https://opsani.com/}}
\end{figure}
 
\subsection{Motivazioni aziendali}

IKS, con il progetto di stage proposto intende principalmente formare personale per l'attività di consulenza al fine di poter fornire ai propri clienti servizi in ambito \textit{Devops}, data la richiesta di mercato sempre maggiore. L'azienda inoltre intendeva valutare l'uso di \gls{container} al fine di ottimizzare l'uso delle risorse dei nodi fisici.

\subsection{Motivazioni personali}

Lo stage curricolare rappresenta la tappa finale del percorso di studi in Informatica presso l'\myUni. Tale attività costituisce un'interessante opportunità per mettere in contatto le realtà aziendali che necessitano di competenze specifiche, con studenti che si trovano ad approcciarsi al mondo del lavoro. A tal proposito ogni anno, tra marzo e aprile l'università, con il patrocinio di Confindustria Padova organizza un apposito evento denominato \textit{STAGE-IT}, per mettere in comunicazione aziende e studenti.\\

Durante \textit{STAGE-IT 2017} ho avuto modo di sostenere sedici colloqui preliminari con varie aziende; colloqui ai quali successivamente ne sono seguiti altri che mi hanno fornito una visione più chiara dei progetti proposti al fine di una scelta finale. Tra le varie proposte, ho stabilito dei criteri così riassunti: personali, professionali, economici e logistici.

\subsubsection{Motivazioni personali}

Le motivazioni personali rappresentano le tematiche di mio interesse, nello specifico cercavo uno stage con le seguenti caratteristiche:

\begin{itemize}
    \item \textbf{Forte orientamento all'open source}. L'uso di tecnologie open source e software libero è in costante crescita tra le aziende informatiche, inoltre avere a disposizione il codice sorgente delle applicazioni che si andranno ad utilizzare permette da una parte di capire il loro funzionamento e dall'altra di potere adattarle secondo i propri bisogni. 
    \item \textbf{Ambito innovativo}. Cercavo uno stage che mi permettesse di rapportarmi con i temi caldi attualmente presenti nel mercato dell'IT, quali ad esempio: \textit{\gls{machine-learning}}, \gls{iot}, \textit{devops}, \gls{blockchain} o \gls{cloud}.
    \item \textbf{Ambito prevalentemente sistemistico}. Il nostro corso di studi si focalizza principalmente nell'ambito sviluppo, l'ambito sistemistico rappresentava una sfida e un'opportunità per integrare le mie conoscenze.
\end{itemize}

\subsubsection{Motivazioni professionali}

\begin{itemize}
    \item \textbf{Ambito innovativo}. Oltre all'interesse personale, approcciare tecnologie innovative pensoi permetta di avere più opportunità lavorative e di confrontarsi con sfide sempre nuove. 
\end{itemize}

\subsubsection{Motivazioni economiche e logistiche}

Alle motivazioni economiche e logistiche, ho preferito anteporre quelle personali e professionali. 

\begin{itemize}
    \item \textbf{Luogo di lavoro}. Sicuramente uno stage vicino casa sarebbe stato comodo ma non mi avrebbe permesso di crescere. La zona industriale di Padova, parecchio lontana sia fisicamente che idealmente alla mia ``comfort-zone'' mi ha sicuramente aiutato in questo.
    \item \textbf{Logistica - Trasporti}. Idealmente avrei voluto minimizzare il tempo di trasporto viaggiando in auto, ma questo non è stato possibile. Il dover coordinare diverse tipologie di trasporto mi ha comunque permesso di migliorare le mie doti organizzative e gestionali. 
    \item \textbf{Rimborsi}. Cercavo un'azienda che fornisse un minimo di rimborso spese. IKS in questo senso ha messo a disposizione dei buoni pasto a parziale copertura delle spese di ristorazione.
\end{itemize}

%**************************************************************
\section{Vincoli}
%**************************************************************
\subsection{Vincoli metodologici}

All'inizio del progetto di stage è stato concordato con il tutor aziendale un piano di lavoro di massima, comprendente la pianificazione dello stage da svolgersi prevalentemente presso la sede dell'azienda ospitante in Corso Stati Uniti 14 bis a Padova.

Regolarmente, (circa una volta la settimana) si sono tenuti incontri diretti con il tutor aziendale  per verificare lo stato di avanzamento del progetto, ridefinire gli obiettivi, affinare
la ricerca e aggiornare il piano stesso di lavoro; il tutor aziendale ha inoltre richiesto un breve report giornaliero delle attività svolte da inviare via mail in modo da essere costantemente aggiornato sul procedere del lavoro e poter tempestivamente chiarire dubbi e fornire chiarimenti.

\subsection{Vincoli temporali}
Lo stage curriculare, all'interno del corso di informatica si deve svolgere in un margine di ore compreso tra le 300 e le 320 ore presso l'azienda ospitante. Il progetto che ho svolto presso IKS ha avuto una durata di 304 ore circa, più alcune ore dedicate da remoto per la stesura di documentazione tecnica. \\

La pianificazione iniziale è stata seguita in linea di massima, ma con l'emergere di altre priorità come lo studio dei protocolli di comunicazione tra \gls{master} e \gls{slave}, inizialmente non previste dal piano di lavoro ha sottratto tempo ad altre attività come lo sviluppo di un \gls{plugin} base da usare come modello di sviluppo.

\subsection{Vincoli tecnologici}

Gli unici vincoli tecnologici imposti dall'azienda circa la piattaforma di \gls{continuous building} da implementare riguardavano il sistema operativo da utilizzare, \gls{CentOS} 7, Jenkins come fulcro centrale del progetto al fine di automatizzare i processi di \gls{build} e Docker come standard ``de-facto'' nelle tecnologie di containerizzazione.

%**************************************************************
\section{Obiettivi}
%**************************************************************

Vengono qui riportati gli obiettivi iniziali dello stage, in forma tabellare.

Ogni obiettivo è descritto da tre attributi:

\begin{itemize}
    \item \textbf{ID}. Un codice identificativo univoco, al fine di tracciare l'obiettivo. Esso è determinato da una lettera indicante la priorità (O, D o F) sequita da un numero progressivo a due cifre.
    \item \textbf{Priorità}. L'importanza attribuita all'obiettivo. Essa può essere:
    \begin{itemize}
        \item \textbf{Obbligatorio}. Obiettivi vincolanti in quanto esplicitamente indicati come primari dall'azienda proponente; 
        \item \textbf{Desiderabile}. Obiettivi non vincolanti o strettamente necessari, ma dal riconoscibile valore
        aggiunto;
        \item \textbf{Facoltativo}. Obiettivi rappresentanti valore aggiunto non strettamente competitivo, giudicati dall'azienda come opzionali.
    \end{itemize} 
    \item \textbf{Descrizione}. Una breve descrizione riepilogativa dell'obiettivo. 
\end{itemize}

\taburowcolors[2] 2{tableLineOne .. tableLineTwo}
\tabulinesep = ^3mm_2mm
\begin{longtabu} to \textwidth {ccX}
    \caption[Obiettivi dello stage]{Obiettivi dello stage}
    \label{tabella:obiettivi-iniziali}
    \endlastfoot
    \rowfont{\bfseries\sffamily\leavevmode\color{white}}
    \rowcolor{tableHeaderRed}
    %\hline
    \textbf{ID} & \textbf{Importanza} & \textbf{Descrizione} \\
    %\hline
    O01 & Obbligatorio & Progettazione dell’architettura di un sistema di \gls{continuous building} \\ %\hline
    O02 & Obbligatorio & Implementazione dell’architettura progettata su sistemi IKS \\ %\hline
    O03 & Obbligatorio & Stesura manuale di installazione \\ %\hline
    O04 & Obbligatorio & Utilizzo della piattaforma implementata per la produzione di \gls{container} Docker \\ %\hline
    O05 & Obbligatorio & Studio di fattibilità sulle modalità di comunicazione \gls{master} e \gls{slave} \\ %\hline
    D01 & Desiderabile & Containerizzazione della piattaforma implementata \\ %\hline
    D02 & Desiderabile & Relazione sui confronti prestazionali tra la piattaforma e la sua controparte containerizzata \\ %\hline
    DO3 & Desiderabile & Individuazione metriche per il monitoraggio della piattaforma \\ %\hline
    DO4 & Desiderabile &  Individuazione procedure per il monitoraggio, backup, ripristino e upgrade dell’infrastruttura \\ %\hline
    DO5 & Desiderabile & Breve relazione che documenti metriche e procedure dei punti DO3-4 \\ %\hline
    F01 & Facoltativo & Implementazione di un plugin base per Jenkins \\ %\hline
\end{longtabu}
\clearpage

In sezione {\hyperref[sec:obiettivi-raggiunti]{``8.1 - Raggiungimento degli obiettivi''}} viene fornito un riepilogo in forma tabellare di tali obiettivi, indicando per ciascuno di essi il grado di soddisfacimento.

%**************************************************************
\section{Pianificazione}
%**************************************************************

In {\hyperref[tab:initial-plan]{tabella 2.2} viene riportata la pianificazione delle attività, con la relativa suddivisione oraria.

\begin{longtabu} to \textwidth{|c|X|}
    \caption[Pianificazione iniziale]{Pianificazione iniziale}
    \label{tab:initial-plan}
    
    \endlastfoot
    \rowfont{\bfseries\sffamily\leavevmode\color{white}}
    \rowcolor{darkblue}
	\textbf{Durata in ore} & \textbf{Descrizione dell'attività} \\\hline
    
    \rowcolor{lightblue}
	\textbf{38} & \textbf{Formazione sulle tecnologie} \\	 
    \hline
    
    \rowcolor{lightblue}
    \textbf{38} & \textbf{Definizione architettura di riferimento e relativa documentazione} \\ \hline
    \multirow{3}{0cm} 
    \textit{12} & 
    \textit{Analisi del problema e del dominio applicativo} \\
    \textit{22} & 
    \textit{Progettazione della piattaforma e relativi test} \\
    \textit{4} & 
    \textit{Stesura documentazione relativa ad analisi e progettazione} \\
    \hline
    
    \rowcolor{lightblue}
    \textbf{38} & \textbf{Implementazione di base} \\ \hline
    \multirow{3}{0cm} 
    \textit{24} & 
    \textit{Implementazione architettura di base} \\
    \textit{10} & 
    \textit{Implementazione test funzionali} \\
    \textit{4} & 
    \textit{Stesura Manuale di installazione} \\
    \hline
    
    \rowcolor{lightblue}
    \textbf{38} & \textbf{Building di container} \\ \hline
    \multirow{2}{0cm} 
    \textit{28} & 
    \textit{Implementazione delle funzionalità di containerizzazione} \\
    \textit{10} & 
    \textit{Testing delle funzionalità} \\
    \hline
    
    \rowcolor{lightblue}
    \textbf{38} & \textbf{Containerizzazione della piattaforma}  \\ \hline
    \multirow{3}{0cm} 
    \textit{25} & 
    \textit{Containerizzazione della piattaforma} \\
    \textit{10} & 
    \textit{Confronti prestazionali} \\
    \textit{3} & 
    \textit{Stesura documentazione relativa alle prestazioni} \\
    \hline
    \newpage
    \rowfont{\bfseries\sffamily\leavevmode\color{white}}
    \rowcolor{darkblue}
    \textbf{Durata in ore} & \textbf{Descrizione dell'attività} \\\hline
    
    \rowcolor{lightblue}
    \textbf{38} & \textbf{Monitoraggio / Backup / Upgrade}  \\ \hline
    \multirow{5}{0cm} 
    \textit{9} & 
    \textit{Individuazione metriche per il monitoraggio} \\
    \textit{15} & 
    \textit{Sviluppo procedure e modalità per il recupero delle metriche} \\
    \textit{5} & 
    \textit{Individuazione procedure di backup e ripristino della piattaforma} \\
    \textit{5} & 
    \textit{Individuazione procedure di upgrade} \\
    \textit{4} & 
    \textit{Stesura documentazione relativa alle metriche e alle procedure individuate} \\
    \hline
    
    \rowcolor{lightblue}
    \textbf{38} & \textbf{Sviluppo plugin sample per Jenkins}  \\ \hline
    \multirow{2}{0cm} 
    \textit{30} & 
    \textit{Implementazione plugin e relativi test funzionali} \\
    \textit{4} &
    \textit{Testing del plugin sviluppato} \\
    \textit{4} & 
    \textit{Stesura documentazione relativa all'uso e manutenzione del plugin} \\
    \hline
    
    \rowcolor{lightblue}
    \textbf{38} & \textbf{Collaudo Finale}  \\ \hline
    \multirow{4}{0cm} 
    \textit{30} & 
    \textit{Collaudo della piattaforma} \\
    \textit{5} & 
    \textit{Stesura documentazione finale} \\
    \textit{1} & 
    \textit{Incontro di presentazione della piattaforma con gli stakeholders} \\
    \textit{2} & 
    \textit{Live demo di tutto il lavoro di stage} \\
    \hline

	\textbf{Totale ore} & \multicolumn{1}{|c|}{\textbf{304}} \\\hline
\end{longtabu}

\newpage

%**************************************************************
\section{Ambiente di Lavoro}
%**************************************************************

\subsection{Metodologia di sviluppo}

IKS adotta la metodologia di sviluppo \gls{agile}, tuttavia, per questo progetto dove la parte di sviluppo è poco prevalente, si è scelto di adottare \gls{cbse}, cioè un processo di sviluppo basato sul riuso di componenti software già esistenti e sulla loro integrazione. \\ 
Visto che le componenti software sono già esistenti, non è detto che esse rispecchino pienamente le funzionalità necessarie al prodotto finale; occorrerà quindi una loro attenta selezione in relazione ai requisiti del prodotto finale e dei bisogni degli \gls{stakeholder}, con i quali andrà instaurato uno stretto rapporto.

\begin{figure}[H]
    \capstart
    \captionsetup{justification=centering}
    \centering
    \includegraphics[width=12cm, keepaspectratio]{../immagini/CBSE.jpg}
    \caption{Rappresentazione dela metodologia CBSE, immagine tratta da \url{http://www.federica.unina.it/ingegneria/ingegneria-software-ii/component-based-software-engineering-sviluppo-cbse/}}
\end{figure}

Il modello a componenti, favorisce la separazione delle responsabilità secondo il \textit{\gls{srp}} e permette di ottenere un basso accoppiamento tra le componenti e il sistema principale.

Jenkins, con la sua architettura a \gls{plugin} ben si presta ad una composizione gerarchica, dove il servizio principale invoca direttamente quelli esposti dai singoli \gls{plugin} che ne estendono le funzionalità.

\subsection{Gestione di progetto}

\subsubsection{Versionamento}
Inizialmente da quanto preventivato, per il versionamento del codice sorgente di alcune applicazioni Maven d'esempio non è stato possibile utilizzare \textit{git}. L'azienda ha però fornito un \gls{repository} \gls{cvs}\footnote{\url{https://savannah.nongnu.org/project/memberlist.php?detailed=1&group=cvs}} che è stato utilizzato principalmente per lo storage e il versionamento di applicazioni Maven e dei \textit{Dockerfile} creati. \\

\gls{cvs} è un sistema di versionamento del codice piuttosto vetusto (non presenta nuove versioni dal 2008) ma essendo uno tra i primi strumenti di versionamento, tra l'altro costituito da software libero, è ancora parecchio usato. \\
Non avendo mai approfondito in dettaglio tale strumento, mi sono avvalso di un client grafico, smartCVS\footnote{\url{https://www.syntevo.com/smartcvs/}}.

\subsection{Documentazione}
La documentazione tecnica è stata redatta utilizzando \textit{Libreoffice 5.4.2}, lo strumento, pur essendo open source non si è rivelato pienamente adatto allo scopo, in quanto non permette una standardizzazione della documentazione (anche se è possibile definire dei template) e della formattazione. \LaTeX\ probabilmente avrebbe rappresentato un'alternativa migliore, in quanto codice pienamente versionabile ma sarebbe di difficile adozione all'interno di un team, in quanto tutti dovrebbero conoscere tale linguaggio. 

\subsection{Ambiente di sviluppo}
Per la piccola parte di sviluppo e personalizzazione di script Groovy si è utilizzato Visual Studio Code. \\
Visual Studio Code\footnote{\url{https://code.visualstudio.com/}}, è un editor open source sviluppato da Microsoft, basato su \textit{Electron}, un framework javascript che permette lo sviluppo di applicazioni desktop.
La caratteristica chiave di Visual Studio Code consiste in un'ampia gamma di \gls{plugin} ed estensioni, che lo rendono un editor versatile e facilmente adattabile a varie esigenze. Esso inoltre risulta configurabile in modo preciso e dettagliato tramite \gls{json}. 

\begin{figure}[H]
    \capstart
    \centering
    \includegraphics[height=3cm, keepaspectratio]{../immagini/VSCode.png}
    \caption{Logo di Visual Studio code}
\end{figure}

%**************************************************************
\section{Analisi preventiva dei rischi}
%**************************************************************
In tabella {\hyperref[tab:ab:rischi]{``2.2 - Analisi dei rischi''}} viene riportata l'analisi dei rischi che ho dovuto affrontare nel corso del progetto. Per ogni rischio viene definito:

\begin{itemize}
    \item Un nome;
    \item Una breve descrizione;
    \item La modalità di rilevamento;
    \item La probabilità di occorrenza;
    \item Il grado di pericolosità;
    \item Il piano di contingenza.
\end{itemize}

\begin{longtabu} to \linewidth { >{\centering}m{.15\linewidth} >{\raggedright}m{.33\linewidth}
>{\raggedright}m{0.25\linewidth} >{\centering}m{0.14\linewidth} }
    \caption[tab-rischi]{Tabella descrittiva dell'analisi dei rischi e attualizzazione}
    \label{tab:rischi}
    \endlastfoot
            
    %\taburowcolors{tableLineOne..tableLineTwo}
    \rowfont{\bfseries\sffamily\leavevmode\color{white}}
    \rowcolor{darkgreen}\textbf{Nome} & \textbf{Descrizione} & \textbf{Rilevamento} & \textbf{Rischio} \\*
    \rowcolor{tableLineOne} \textbf{Scarsa esperienza} &
    {\small L'ambito innovativo e la scarsa conoscenza pregressa potrebbero comportare ritardi nell'esecuzione delle attività previste} & {\small Occorrerà comunicare al tutor aziendale eventuali 
	difficoltà} & \shortstack{Occorrenza: \\ \textbf{Medio-alta} \vspace{0.6em}\\ Pericolosità: \\ \textbf{Alta} }\\*
    \rowcolor{tableLineTwo} Piano di contingenza: & \multicolumn{3}{m{0.7828\linewidth}}{\small Occorre segnalare tempestivamente al tutor aziendale eventuali problematiche e approfondire la formazione} \\*
        
    \rowfont{\bfseries\sffamily\leavevmode\color{white}}
    \rowcolor{darkgreen}\textbf{Nome} & \textbf{Descrizione} & \textbf{Rilevamento} & \textbf{Rischio} \\*
	\rowcolor{tableLineOne} \textbf{Disponibilità temporali} &{\small La disponibilità dei mezzi di trasporto pubblici ed eventuali imprevisti possono causare l'impossibilità nell'essere in sede}&{\small Occorrerà comunicare al tutor aziendale eventuali ritardi}& \shortstack{Occorrenza: \\ \textbf{Medio-bassa} \vspace{0.6em}\\ Pericolosità: \\ \textbf{Medio-alta} }\\*
	\rowcolor{tableLineTwo} Piano di contingenza:&\multicolumn{3}{m{0.7828\linewidth}}{\small Verranno fissate delle giornate di recupero, se il lavoro riguarda la documentazione potrà essere svolto da remoto}\\*
            
    \rowfont{\bfseries\sffamily\leavevmode\color{white}}
    \rowcolor{darkgreen}\textbf{Nome} & \textbf{Descrizione} & \textbf{Rilevamento} & \textbf{Rischio} \\*
	\rowcolor{tableLineOne} \textbf{Tecnologie da usare} &{\small Il tempo di apprendimento per le tecnologie potrebbe causare ritardi nello svolgimento dei lavori}&{\small Andranno evidenziate eventuali lacune al tutor aziendale}& \shortstack{Occorrenza: \\ \textbf{Media} \vspace{0.6em}\\ Pericolosità: \\ \textbf{Medio-alta} }\\*
	\rowcolor{tableLineTwo} Piano di contingenza:&\multicolumn{3}{m{0.7828\linewidth}}{\small Andranno aumentate le ore disponibili in attività di formazione}\\*

    \rowfont{\bfseries\sffamily\leavevmode\color{white}}
    \rowcolor{darkgreen}\textbf{Nome} & \textbf{Descrizione} & \textbf{Rilevamento} & \textbf{Rischio} \\*         
    \rowcolor{tableLineOne} \textbf{Problemi hardware}&{\small Potrebbero verificarsi guasti hardware i quali comporterebbero perdite di dati e/o di tempo}&{\small Andrà segnalato al tutor aziendale ogni eventuale malfunzionamento o guasto}& 
    \shortstack{Occorrenza: \\ \textbf{Medio-bassa} \vspace{0.6em}\\ Pericolosità: \\ \textbf{Bassa} } \\*
    \rowcolor{tableLineTwo} Piano di contingenza:&\multicolumn{3}{m{0.7828\linewidth}}{\small Il team sistemistico di IKS provvederà al ripristino dell'infrastruttura}\\*
    
    \rowfont{\bfseries\sffamily\leavevmode\color{white}}
    \rowcolor{darkgreen}\textbf{Nome} & \textbf{Descrizione} & \textbf{Rilevamento} & \textbf{Rischio} \\*
	\rowcolor{tableLineOne} \textbf{Strumenti software} &{\small Il progetto prevede l'interoperabilità di componenti software che potrebbero rivelarsi incompatibili tra loro}&{\small Andrà specificata per ogni componente la versione, facendo riferimento alla documentazione ufficiale, se presente.}& 
    \shortstack{Occorrenza: \\ \textbf{Medio-alta} \vspace{0.6em}\\ Pericolosità: \\ \textbf{Bassa} }\\*
	\rowcolor{tableLineTwo} Piano di contingenza:&\multicolumn{3}{m{0.7828\linewidth}}{\small Andranno scelte componenti che siano compatibili, nel limite del possibile}\\*
\end{longtabu}
             % Processi
% !TEX encoding = UTF-8
% !TEX TS-program = pdflatex
% !TEX root = ../tesi.tex

%**************************************************************
\pagestyle{IHA-fancy-style}
\chapter{Tecnologie adottate}
\label{cap:tecnologie-adottate}
%**************************************************************

\intro{In questo capitolo vengono esposte brevemente le tecnologie, gli strumenti e i linguaggi utilizzati durante lo svolgimento dello stage. Per ognuno di essi, verrà fornita una panoramica di alto livello, che verrà eventualmente approfondita nei capitoli successivi.}\\

%**************************************************************
\section{Tecnologie software}

\subsection{CentOS 7}

\gls{CentOS} 7 \footnote{\url{https://www.centos.org/}} è la distribuzione \gls{GNU}/\gls{Linux} di riferimento all'interno di IKS. \gls{CentOS} è una distribuzione di tipo aziendale, basata su \gls{Red Hat} Enterprise \gls{Linux} e si distingue da quest'ultima per la mancanza di supporto tecnico fornito da \gls{Red Hat} e per il cambio di logo.\\
\gls{CentOS} fa dell'affidabilità e della robustezza i suoi punti cardine, rendendola una delle migliori distribuzioni \gls{GNU}/\gls{Linux} attualmente presenti sul mercato in ambito server.

\begin{figure}[H]
    \capstart
    \centering
    \includegraphics[height=4cm, keepaspectratio]{../immagini/CentOS-logo.png}
    \caption{Logo di CentOS 7}
\end{figure}

\newpage

\subsection{Bash 4.4.5}

\gls{bash} \footnote{\url{https://www.gnu.org/software/bash/}} è una shell, cioè un interprete a linea di comando frequentemente presente all'interno di sistemi \gls{Linux}. Essa fornisce un semplice linguaggio di scripting nativo che tramite l'impostazione di variabili, comandi e strutture di controllo del flusso permette di eseguire compiti complessi, favorendone l'automazione.

\begin{figure}[H]
    \capstart
    \centering
    \includegraphics[width=5cm, keepaspectratio]{../immagini/bash-logo.png}
    \caption{Logo Bash 4.4.5}
\end{figure}

\subsection{JDK 8 update 151}

Il \gls{jdk} è una raccolta di strumenti essenziali per la programmazione in linguaggio Java. Esso fornisce, tra le altre, una componente per la compilazione del codice sorgente, un interprete per il \gls{bytecode} generato dal compilatore, uno strumento per la generazione automatica di documentazione a partire da particolari commenti presenti nel codice sorgente, un \textit{\gls{debugger}} e molto altro.

\begin{figure}[H]
    \capstart
    \centering
    \includegraphics[width=5cm, keepaspectratio]{../immagini/JDK.jpg}
    \caption{Logo JDK}
\end{figure}

\subsection{Apache Tomcat 9.0.1}

Apache Tomcat \footnote{\url{https://tomcat.apache.org/}} è un \textit{\gls{application-server}} una piattaforma software per l'esecuzione di applicazioni Web sviluppate in linguaggio Java. Esso permette una facile gestione di applicazioni web, consentendo di definire varie politiche di gestione degli utenti e una dettagliata configurazione di porte e altre impostazioni di rete.

\begin{figure}[H]
    \capstart
    \centering
    \includegraphics[width=5cm, keepaspectratio]{../immagini/Tomcat-logo.png}
    \caption{Logo Apache Tomcat 9.0.1}
\end{figure}

\subsection{Jenkins 2.73.3}

Jenkins \footnote{\url{https://jenkins.io/}} è uno strumento di \gls{integrazione-continua} open source e rappresenta uno dei principali strumenti di automazione di compiti. Esso permette agli utenti, tramite la definizione di \gls{project} (in passato definiti jobs) di automatizzare una serie complessa di compiti, quali ad esempio la compilazione di una componente software a partire dal suo codice sorgente, l'esecuzione dei test e la visualizzazione dei loro risultati e il \gls{deploy} dell'\gls{artefatto} ottenuto verso un \glossary{repository} dei prodotti.\\

Jenkins è facilmente installabile e configurabile ed è funzionante ``out of the box'' senza particolari impostazioni; presenta un ricco ecosistema di \gls{plugin} messi a disposizione dalla comunità ed è estensibile, tramite questi ultimi, qualora le funzionalità desiderate non fossero già disponibili.\\

Jenkins si basa su un'architettura \gls{master} - \gls{slave}, dove un'istanza principale (il \gls{master}) coordina i vari compiti tra uno o più \gls{slave}, i quali rappresentano dei meri esecutori dei lavori a loro assegnati.

\begin{figure}[H]
    \capstart
    \centering
    \includegraphics[height=3.5cm, keepaspectratio]{../immagini/jenkins.png}
    \caption{Logo Jenkins 2.73.3}
\end{figure}

\subsection{Docker 17.10 Community Edition}

Docker \footnote{\url{https://www.docker.com/}} è un progetto open source che consente la creazione, l'istanziazione e l'esecuzione di processi isolati, definiti \gls{container} in un sistema ospite. \\

L'idea alla base di Docker è la containerizzazione, una particolare forma di virtualizzazione che basandosi sui \textit{cgroups}, una funzionalità del kernel \gls{Linux} che rappresenta dei gruppi di controllo in grado di offrire applicazioni complete viste come risorse isolate, senza che sia necessario istanziare delle macchine virtuali pienamente operative. \\

Possiamo immaginare un container in ambito informatico, come la sua controparte legata ai trasporti fisici di merci. In un container che siamo abituati a vedere trasportato da un treno o stivato in una nave, la merce viene racchiusa al suo interno e spedita in luoghi specifici. Il concetto si ritrova in campo informatico, dove in un container è possibile includere un'applicazione software con tutte le sue dipendenze e con le utilità necessarie al suo corretto funzionamento, in modo isolato dal resto del sistema. 

Questo permette di avere applicazioni portabili indipendentemente dal sistema operativo sul quale andranno in esecuzione, un uso efficiente delle risorse di sistema, la scalabilità delle applicazioni e la garanzia che ogni \gls{container} sarà un processo isolato dagli altri \gls{container} e dal sistema ospite stesso, consentendo quindi ogni volta che viene mandato in esecuzione d'avere un ambiente ``pulito'', con le caratteristiche desiderate. 

\begin{figure}[H]
    \capstart
    \centering
    \includegraphics[width=5cm, keepaspectratio]{../immagini/docker.png}
    \caption{Logo Docker 17.10}
\end{figure}

\subsection{Apache Maven 3.5.2}

Apache Maven \footnote{\url{https://maven.apache.org/}} è uno strumento di automazione delle \gls{build}, usato principalmente per progetti \textit{Java}. Grazie a Maven, è possibile gestire le dipendenze di un determinato progetto e definirne particolari aspetti, quali ad esempio le politiche di \textit{\gls{deploy}} e molto altro. \\

\begin{figure}[H]
    \capstart
    \centering
    \includegraphics[width=5cm, keepaspectratio]{../immagini/maven-logo.png}
    \caption{Logo Apache Maven 3.5.2}
\end{figure}

Concetto fondamentale di Maven è il \textit{\gls{pom}\footnote{\url{https://maven.apache.org/pom.html}}}, un particolare file di configurazione, basato su \gls{xml} che descrive il progetto, con le sue dipendenze e la sua configurazione. Ogni progetto Maven dovrà essere dotato di almeno un \gls{pom}, e rispettare la struttura in esso definita.

\begin{figure}[H]
    \capstart
    \centering
    \includegraphics[width=12cm, keepaspectratio]{../immagini/maven-directory.png}
    \caption[Struttura di un progetto Maven, così come definita da un generico POM]{Struttura di un progetto Maven, così come definita da un generico POM, immagine tratta da \url{http://www.murraywilliams.com/2012/04/maven-and-jpa-programming/}}
\end{figure}

\subsection{cadvisor 0.28.2}

cadvisor \footnote{\url{https://github.com/google/cadvisor}} è uno strumento sviluppato da Google per il monitoraggio e l'analisi in real-time delle performance e del consumo di risorse di \gls{container}. \\
Fornisce inoltre supporto nativo ai \gls{container} Docker.

\begin{figure}[H]
    \capstart
    \centering
    \includegraphics[height=4cm, keepaspectratio]{../immagini/cadvisor.png}
    \caption{Logo cadvisor 0.28.2}
\end{figure}

%**************************************************************
\section{Linguaggi di Programmazione}
%**************************************************************
\subsection{Java}

Java è un linguaggio orientato agli oggetti, fortemente tipizzato e focalizzato sull'essere multi-piattaforma. Tale linguaggio grazie alla \gls{jvm} permette una notevole portabilità del codice sorgente, che verrà compilato sotto forma di \gls{bytecode}. \\

Java è inoltre il principale linguaggio col quale vengono sviluppati \gls{plugin} per Jenkins, essendo esso stesso scritto prevalentemente in tale linguaggio.

\begin{figure}[H]
    \capstart
    \centering
    \includegraphics[height=4cm, keepaspectratio]{../immagini/java-logo-vector.png}
    \caption{Logo Java 8}
\end{figure}

\newpage

\subsection{Groovy}

Groovy \footnote{\url{http://groovy-lang.org/}} è un linguaggio per la piattaforma Java ma indipendente da esso, pur avendo una sintassi molto simile. \`E sviluppato dall'Apache Software Foundation con l'intento di essere facile da imparare, semplice da integrare con librerie e componenti scritte in Java, con funzionalità aggiuntive e un ricco ecosistema di framework a supporto. Può inoltre essere visto come un linguaggio di scripting semplice e maneggevole per la scrittura di test concisi e manutenibili oltre che per l'automazione di vari compiti. 

\begin{figure}[H]
    \capstart
    \centering
    \includegraphics[width=5cm, keepaspectratio]{../immagini/groovy.png}
    \caption{Logo Groovy}
\end{figure}             % Kick-Off
% !TEX encoding = UTF-8
% !TEX TS-program = pdflatex
% !TEX root = ../tesi.tex

%**************************************************************
\pagestyle{IHA-fancy-style}
\chapter{Realizzazione del sistema}
\label{cap:realizzazione-del-sistema}
%**************************************************************

\intro{In questo capitolo vengono esposte le attività relative alla realizzazione del sistema; dalla fase di analisi e progettazione, all'effettiva implementazione. Viene inoltre riportato una sintesi dello studio dei protocolli e dei metodi di comunicazione tra \gls{master} e \gls{slave} riferito ad ambienti containerizzati.}\\

%**************************************************************
\section{Descrizione generale dell'architettura individuata}
%**************************************************************

In questa sezione si intende analizzare i limiti dell'architettura attuale e si riportano brevemente i vantaggi della soluzione implementata.

\subsection{Limiti dell'architettura attuale}
\label{subsec:limiti}

Tipicamente si ha un'istanza \gls{master} di Jenkins ospitata su un server fisico o virtuale, la quale coordina il lavoro e dirige lo svolgimento dello stesso, delegando più compiti a vari nodi \gls{slave}. Ciascuno slave a sua volta, sarà ospitato da un \gls{host} fisico o virtuale, il quale dovrà provvedere le risorse adeguate in termini di memoria, capacità di elaborazione, disponibilità nell'esecuzione dei compiti eccetera. Tale architettura risulta essere poco flessibile e rappresenta un \textit{Single Point of Failure}; qualora il nodo \gls{slave} risulti offline o irraggiungibile per qualche motivo, il nodo \gls{master} non sarà in grado di mandare in esecuzione su di esso alcun \gls{project}. I nodi \gls{slave} inoltre andranno preparati in precedenza, configurando manualmente la macchina \gls{host}, installandovi eventuali dipendenze, quali ad esempio: 

\begin{itemize}
    \item Strumenti di automazione delle \gls{build} come Maven, Gradle o Ant; 
    \item Framework specifici per le applicazioni sviluppate;
    \item Utilità di compilazione e supporto allo sviluppo come il \gls{jdk}.
\end{itemize}

\`E impensabile dover ripetere manualmente tali operazioni per ogni \gls{host}; senza contare che potrebbe essere necessario avere più utilità, anche in conflitto tra loro (come ad esempio diverse versioni del \gls{jdk}, al fine di garantire la retro compatibilità delle applicazioni sviluppate). Tale architettura è difficilmente manutenibile, troppo rigida e quasi impossibile da adattare alle esigenze; ne consegue che ogni nodo slave una volta configurato sarà adatto solo per determinati tipi di \gls{build}.\\

L'architettura presenta inoltre un uso inefficiente delle risorse. Se l'intera macchina è dedicata univocamente a svolgere il ruolo di \gls{slave}, si troverà ad eseguire i compiti assegnati dal \gls{master} per un certo tempo, ma la maggior parte lo trascorrerà in uno stato di inattività. \\

In \hyperref[fig:impl1]{figura 4.1} è illustrata la configurazione iniziale che ho utilizzato presso IKS per l'attività di formazione e di sperimentazione. Tale infrastruttura rispecchia pienamente i limiti appena descritti, in questo caso il ruolo di \gls{master} viene assunto da un'istanza di Jenkins presente su un portatile con Windows 10, mentre il ruolo di \gls{slave} è ricoperto dalla macchina virtuale \gls{Linux} remota; in questa configurazione, si sono ottenuti come artefatti, sia applicazioni generate da progetti Maven (in formato ``.\gls{war}'') che immagini Docker. Per raggiungere tali risultati è stato necessario installare e configurare nello slave remoto sia Maven che Docker.

\begin{figure}[H]
    \capstart
    \centering
    \includegraphics[width=12cm, keepaspectratio]{../immagini/impl1.png}
    \caption{Visione d'alto livello dell'architettura iniziale}
    \label{fig:impl1}
\end{figure}

%**************************************************************
\subsection{Vantaggi della soluzione}
%**************************************************************

Le tecnologie di containerizzazione fornite da Docker sono risultate essere la scelta opportuna per sopperire i limiti dell'architettura esposti in \hyperref[subsec:limiti]{sezione 4.1.1}. I \gls{container} possono essere messi in esecuzione con un solo comando da terminale e non necessitano di grosse configurazioni; forniscono inoltre ogni qualvolta vengano eseguiti, un'ambiente ``nuovo'' (non contaminato da installazioni o configurazioni precedenti) e isolato dal resto del sistema ospite. Ecco quindi che un nodo \gls{host} non dovrà più essere dedicato solamente al ruolo di \gls{slave}, ma quando esso riceverà dei compiti da eseguire dal \gls{master} potrà allocare uno o più \gls{container} che svolgano il lavoro e de-allocarli una volta che abbiano raggiunto lo scopo, dedicando nel tempo rimanente risorse ad altre attività.\\

La \hyperref[fig:impl5]{figura 4.2} rappresenta la situazione ideale, con un \gls{master} centrale ospitato su una macchina dedicata, il quale coordina e dirige vari \gls{project} su un numero imprecisato di slave allocati su sua richiesta. Ogni nodo non assumerà più quindi il ruolo di \gls{slave} ma piuttosto quello di fornitore di \gls{slave}, potendone allocare diversi secondo le risorse a sua disposizione. \`E stata implementata, un'architettura ancor più evoluta, dove anche l'istanza \gls{master} di Jenkins è contenuta in un \gls{container}; questo permettere, ad esempio, di poter utilizzare l'intero sistema su una sola macchina, riducendo al minimo la configurazione e le attività di mantenimento dell'infrastruttura.

\begin{figure}[H]
    \capstart
    \centering
    \includegraphics[width=7cm, keepaspectratio]{../immagini/impl5.png}
    \caption{Rappresentazione di alto livello della soluzione individuata}
    \label{fig:impl5}
\end{figure}

%**************************************************************
\section{Implementazione}
%**************************************************************

In un primo momento ho proceduto con la configurazione del sistema, come riportato in \hyperref[fig:impl1]{figura 4.1}, in modo da comprenderne il funzionamento e studiarne le peculiarità, questo mi ha permesso di ampliare notevolmente le mie conoscenze nell'ambito di amministrazione di sistema utilizzando varie tecnologie proprie del mondo \gls{Linux}, come ad esempio \gls{systemd}. 

Parte consistente dell'implementazione del sistema è stata la progettazione e la costruzione delle immagini dei \gls{container}. Gli sviluppatori di Jenkins mettono già a disposizione delle immagini sia per i \gls{master} che per gli \gls{slave}, tuttavia l'azienda ha preferito la creazione da zero in modo da avere delle immagini pienamente adatte alle proprie esigenze. In \hyperref[fig:images]{figura 4.3} è rappresentata la struttura gerarchica delle immagini, così come definite all'interno dei \textit{Dockerfile}\cite{site:dockerfile-reference}.
\begin{figure}[H]
    \capstart
    \centering
    \includegraphics[width=12cm, keepaspectratio]{../immagini/structure.png}
    \caption{Visione astratta della struttura del sistema}
    \label{fig:images}
\end{figure}

A partire dall'immagine pre-costruita di \gls{CentOS} 7, ho definito una serie di \textit{Dockerfile}, arrivando a definire tutta l'architettura del sistema. Nel listing 4.1 viene riportato il \textit{Dockerfile} dell'immagine di Java 8 update 151, seguito da un breve commento.

\lstinputlisting[language=docker, firstline=1, lastline=14, caption={Esempio di \textit{Dockerfile}}]{./code/dockerfile.txt}  

In riga 1 di Listing 4.1, con il comando ``FROM'', verrà utilizzata l'immagine ufficiale di \gls{CentOS} 7 come base dalla quale derivare. \\
In riga 3 con il comando ``MAINTAINER'' viene riportato il creatore dell'immagine ed un suo contatto. \\
Nelle righe 6-9 tramite il comando ``ENV'' vengono definite delle variabili d'ambiente. \\
In riga 12 l'archivio contenente i binari di installazione di Java viene decompresso nella directory \textit{/opt} con il comando ``ADD''. \\
Infine, con il comando ``CMD'' viene eseguito un comando \textit{echo} di conferma che l'installazione è andata a buon fine. In questo \textit{Dockerfile} sono definiti otto comandi, quindi l'immagine risultante avrà otto strati.\\

Un'\gls{immagine} consiste in un insieme di strati, ciascuno dei quali porta a delle modifiche del filesystem del \gls{container} che andrà in esecuzione, secondo determinati parametri. Le immagini possono derivare gerarchicamente l'una dall'altra in un modo molto simile all'estensione di classi in un linguaggio di programmazione orientato agli oggetti. \\ 

Per definire un'immagine, occorre riassumere l'insieme di comandi\cite{site:docker-command-for-cli} in un \textit{Dockerfile}; semplificando, quando l'immagine verrà costruita con un comando di build, ogni comando del \textit{Dockerfile} diventerà uno strato dell'immagine andando ad aumentarne la dimensione; al fine di ridurre tali dimensioni\cite{site:docker-development-best-practice} esistono delle best practice per la stesura dei \textit{Dockerfile}\cite{site:best-practice-for-dockerfile}.

Durante lo stage è stata data particolare attenzione all'ottimizzazione delle immagini, in particolare alla riduzione della loro dimensioni, per meglio comprendere quanto la dimensione possa variare, viene riportato di seguito un confronto tra comandi simili, presi direttamente dal \textit{Dockerfile} dell'immagine di Apache Tomcat 9.0.1.

\lstinputlisting[language=docker, firstline=1, lastline=8, caption={Esempio di comando RUN in un \textit{Dockerfile}}]{./code/docker-example.txt}  

Nelle righe 1-3 di Listing 4.2 vengono definite tramite il comando ``ENV'' 3 variabili d'ambiente con i rispettivi valori. Esse rappresentano, in ordine:
\begin{itemize}
    \item il percorso d'installazione di Tomcat;
    \item la major version di Tomcat;
    \item la versione specifica di Tomcat.
\end{itemize} 
Successivamente in riga 5 con il comando \textit{wget} viene recuperato l'archivio d'installazione di Tomcat dal sito ufficiale; tale archivio viene decompresso con il comando tar e successivamente eliminato ed infine la directory con i file estratti viene trasferita nel percorso d'installazione scelto. Si noti come le quattro operazioni siano incluse in un unico comando ``RUN'' invece che in quattro; questo serve a ridurre il numero di strati (e quindi la dimensione finale) dell'immagine.\\

L'installazione può avvenire anche se disponiamo già dell'archivio senza doverlo scaricare dalla rete (questo può tornare utile in caso la connessione non sia disponibile), nel listing 4.3 viene riportato lo stesso processo di installazione facendo uso del comando ``COPY''; tale comando andrà a copiare l'archivio all'interno del container.

\lstinputlisting[language=docker, firstline=10, lastline=14, caption={Esempio di installazione offline di Tomcat all'interno di un \textit{Dockerfile}}]{./code/docker-example.txt}  

\`E possibile fare ancora meglio, utilizzando il comando ``ADD''. Tale comando è in grado di scompattare direttamente un'archivio ``.tar'' o ``.tar.gz'' in un percorso scelto, favorendo la leggibilità del \textit{Dockerfile} e minimizzando il numero di strati prodotti a uno. Un esempio d'uso di ``ADD'' è riportato nel listing 4.4, in una sola riga l'archivio viene decompresso nella directory \textit{/opt}.

\lstinputlisting[language=docker, firstline=16, lastline=16, caption={Uso del comando ``ADD'' in un \textit{Dockerfile}}]{./code/docker-example.txt}  

Nella \hyperref[tab:dim]{tabella 4.1} viene riportato un confronto tra le dimensioni delle varie immagini utilizzando i diversi tipi di comandi. Il risparmio calcolato in termini di megabyte oscilla tra i 200 e i circa 300 megabyte.

\taburowcolors[2] 2{tableLineOne .. tableLineTwo}
\tabulinesep = ^3mm_2mm
\begin{longtabu} to \textwidth {ccc}
    \caption[Dimensioni delle immagini]{Dimensioni delle immagini}
    \label{tab:dim}
    \endlastfoot
    \rowfont{\bfseries\sffamily\leavevmode\color{white}}
    \rowcolor{mymauve}
    %\hline
    \textbf{Nome dell'immagine} & \textbf{MB usando RUN} & \textbf{MB usando ADD} \\
    %\hline
    java8u151 & 771 & 581\\ %\hline
    tomcat9.0.1 & 807 & 608\\ %\hline
    jenkins-master & 926 & 727\\ %\hline
    jenkins-ssh-slave & 1260 & 955\\ %\hline
   
\end{longtabu}

Una volta create le immagini, ho proceduto alla configurazione dei \gls{container} andando a definire un volume\footnote{\url{https://docs.docker.com/engine/admin/volumes/volumes/}} per poter gestire agevolmente la directory principale di Jenkins, in modo da semplificare le operazioni legate all'esercibilità, come descritto nel \hyperref[cap:esercibilità]{quinto capitolo}. Sono inoltre stati definiti e adattati degli script Groovy al fine di automatizzare alcune operazioni di configurazione del Jenkins \gls{master}; infine ho provveduto ad una selezione di \gls{plugin} per la fornitura di \gls{container} e monitoraggio. 

\subsection{Difficoltà riscontrate e punti di attenzione}

Le principali difficoltà riscontrate riguardano:

\begin{itemize}
    \item \textbf{Documentazione obsoleta o sparsa}. La documentazione di Jenkins per molti aspetti è letteralmente un ``work in progress'' e quindi per buona parte completamente mancante od obsoleta. Essa non presenta una chiara strutturazione ma è sparsa in più collegamenti di difficile reperimento;
    \item \textbf{Non univocità nella configurazione di un proxy}. Per accedere alla rete esterna, spesso è comune l'uso di un proxy in ambito aziendale. La sua configurazione è stata un aspetto non banale in quanto è possibile impostarlo in un'unica maniera ma vanno attuate circa sei / sette impostazioni distinte, a livello di sistema operativo, a livello di Docker, delle immagini, dei container;
    \item \textit{Conflitti di permessi}\cite{site:understanding-permissions-in-docker}. Soprattutto nel caso dei backup dei volumi\cite{site:understanding-docker-volumes} si sono riscontrati comportamenti curiosi e difficilmente interpretabili con la gestione dei permessi dei file in \gls{Linux}. 
\end{itemize}

\newpage

%**************************************************************
\section{Modalità di comunicazione master/slave}
%**************************************************************

Questa sezione riporta una sintesi dello studio di fattibilità inerente le modalità di comunicazione tra l'istanza \gls{master} di Jenkins e gli \gls{slave} \gls{container} che verranno allocati su richiesta del master\cite{site:jenkins-distributed-builds}. Verranno introdotti brevemente i protocolli utilizzati illustrando il loro funzionamento contestualizzato nel dominio applicativo del progetto. \\

\textit{Yet Another Docker Plugin}\footnote{\url{https://plugins.jenkins.io/yet-another-docker-plugin}}, il \gls{plugin} selezionato per la fornitura e la messa in esecuzione degli slave container, supporta due modalità di comunicazione: \gls{ssh} e \gls{jnlp}.  

\subsection{SSH}
\label{subsec:ssh}

Il protocollo SSH permette di stabilire connessioni di rete sicure, tramite uno scambio di chiavi cifrate tra un server e un client, come illustrato in \hyperref[fig:ssh-connection]{figura 4.4}.

\begin{figure}[H]
    \capstart
    \label{fig:ssh-connection}
    \captionsetup{justification=centering}
    \centering
    \includegraphics[width=9cm, keepaspectratio]{../immagini/ssh-auth.png}
    \caption[Esempio di creazione della connessione tramite SSH]{Esempio di creazione della connessione tramite SSH, immagine tratta da \\ \url{https://www.ssh.com/ssh/protocol/}}
\end{figure}

Questo protocollo di rete è risultato essere la via preferibile e la più semplice per stabilire una connessione sicura tra master e slave in quanto Jenkins contiene un'implementazione del client SSH che può essere utilizzata per stabilire connessioni con gli slave in modo semi automatico. \\
Un generico slave container\cite{site:dockerize-a-ssh-service}, dovrà contenere nella directory \textit{\textasciitilde/.ssh/authorized\_keys} la chiave pubblica del master generata in precedenza; per agevolare tale processo l'immagine base ssh da me definita va già a caricare nel container la chiave pubblica, impostando gli opportuni permessi di sicurezza. A questo punto basterà solamente configurare in Jenkins le credenziali, impostando la chiave privata e lo \textit{Yet Another Docker Plugin} svolgerà tutte le operazioni in modo completamente automatizzato.

\subsection{JNLP}
\label{subsec:jnlp}

\begin{shadequote}
    The Java Network Launch Protocol (\gls{jnlp}) enables an application to be launched on a client desktop by using resources that are hosted on a remote web server..\par\emph{Oracle about JNLP\footnote{\url{https://docs.oracle.com/javase/tutorial/deployment/deploymentInDepth/jnlp.html}}}
\end{shadequote}

Il protocollo \gls{jnlp}, tramite dei descrittori in \gls{xml} specifica le risorse da recuperare e le modalità di esecuzione di applicazioni \gls{jws}, un particolare \textit{framework} Java che permette agli utenti di avviare applicazioni dalla rete direttamente tramite un \textit{browser}. \\
JNLP permette di scaricare localmente tutte le risorse di un'applicazione in modo da eseguirla rapidamente, andando a verificare periodicamente la presenza di aggiornamenti. \\
Fino a luglio 2017, \textit{Yet Another Docker Plugin} non forniva pieno supporto a tale protocollo.
All'interno dell'architettura implementata JNLP si è rivelato particolarmente difficile da configurare, a causa dei limiti imposti dal proxy aziendale\cite{site:troubleshoot-jnlp}; per sfruttare appieno questo protocollo si deve infatti configurare opportunamente il proxy in più punti:

\begin{itemize}
    \item A livello di sistema operativo della macchina \gls{host};
    \item Nel \gls{demone} Docker;
    \item All'interno del \gls{container} master;
    \item All'interno del container slave;
    \item Nelle opzioni di Java, passate tramite comando tramite la variabile JAVA\_OPTS;
    \item Nel file \textit{settings.xml} di Maven\cite{site:maven-settings-reference}.
\end{itemize}

Tutte queste impostazioni sono difficilmente automatizzabili soprattutto confrontandole con il protocollo \gls{ssh} descritto in \hyperref[subsec:ssh]{sezione 4.3.1}.

\begin{figure}[H]
    \capstart
    \label{fig:jnlp-connection}
    \captionsetup{justification=centering}
    \centering
    \includegraphics[width=13cm, keepaspectratio]{../immagini/JNLPsd.png}
    \caption{Diagramma di sequenza della connessione via JNLP}
\end{figure}

In \hyperref[fig:jnlp-connection]{figura 4.5}, tramite un diagramma di sequenza \gls{umlg} viene rappresentata la procedura di connessione tra master e slave. Il processo è completamente automatizzabile e a carico del \textit{Yet Another Docker Plugin} ma non è detto che questo vada a buon fine, in quanto il \textit{download} dei ``.\gls{jar}'' potrebbe essere bloccato dal proxy, se presente.               % Concept Preview
% !TEX encoding = UTF-8
% !TEX TS-program = pdflatex
% !TEX root = ../tesi.tex

%**************************************************************
\pagestyle{IHA-fancy-style}
\chapter{Esercibilità}
\label{cap:esercibilità}
%**************************************************************

\intro{Questo capitolo ha lo scopo di approfondire gli aspetti legati all'esercibilità della piattaforma implementata. Verranno quindi elencati i plugin scelti per Jenkins e le procedure di monitoraggio, upgrade, backup e ripristino individuate.} 

%**************************************************************
\section{Monitoraggio}
%**************************************************************

Il monitoraggio costituisce un'attività essenziale nel valutare le performance del sistema, prevenirne i guasti ed attuare tempestivamente le necessarie misure correttive. Nell'architettura individuata le macro aree da monitorare sono essenzialmente due: i \gls{container} e il sistema containerizzato (Jenkins). \\
Per i \gls{container} in esecuzione è interessante conoscere le risorse utilizzate, quali ad esempio: la quantità di memoria, la percentuale di utilizzo del processore e l'uso della rete e come questi parametri cambino in base alle attività del sistema containerizzato. Lo strumento che ho reputato più adatto a tale scopo è \hyperref[subsec:cadvisor]{cadvisor} di Google, il quale permette un monitoraggio, tramite grafici e statistiche in real-time di tutti questi aspetti, purtroppo non gestisce direttamente la persistenza dei log e dei dati raccolti, ma è integrabile con altri strumenti dedicati\footnote{Una panoramica di tali strumenti è reperibile nella documentazione del repository di cadvisor, al link: \url{https://github.com/google/cadvisor/blob/master/docs/storage/README.md}}. 

\begin{figure}[H]
    \capstart
    \centering
    \captionsetup{justification=centering}
    \includegraphics[width=12cm, keepaspectratio]{../immagini/memory.png}
    \caption[Uso della memoria da parte di un container Jenkins \gls{master}]{Uso della memoria da parte di un container Jenkins \gls{master}, si noti il decadimento dovuto alla conclusione di un \gls{project}}
\end{figure}
 
Per il monitoraggio di Jenkins esistono già svariate soluzioni tramite \gls{plugin} o l'adozione di strumenti specifici, quali ad esempio Datadog\footnote{\url{https://www.datadoghq.com/}}. Datadog sarebbe stata la mia prima scelta, vista la completezza dei dati raccolti e l'agevole visualizzazione degli stessi in comode e dettagliate \textit{dashboard}, tuttavia è un software a pagamento (disponibile gratuitamente per un periodo di prova di soli 14 giorni). Ho quindi ripiegato sulla selezione di \gls{plugin} per Jenkins.

\subsection{Scelta dei plugin per Jenkins}
\label{subsec:scelta-plugin}

Dopo aver valutato diverse opzioni, la scelta è ricaduta su quattro \gls{plugin}, nello specifico:

\begin{itemize}
    \item \textbf{Metrics Plugin\footnote{\url{https://wiki.jenkins.io/display/JENKINS/Metrics+Plugin}}}. Espone varie \gls{apig}, tramite \gls{json} in modo da poterle poi riusare
    con altri sistemi. Particolarmente significative sono le metriche sullo stato dei \gls{plugin}
    caricati, sul tempo impiegato dai \gls{project}, sullo stato delle \gls{build} e sul numero di build eseguite con successo o fallite per project;
    \item \textbf{Monitoring Plugin\footnote{\url{https://wiki.jenkins.io/display/JENKINS/Monitoring}}}.Permette di monitorare tramite
    JavaMelody\footnote{Libreria per il monitoraggio di applicazioni \gls{Java EE} \url{https://github.com/javamelody/javamelody/}} vari aspetti di Jenkins tra cui
    sicurezza, tempo di attività dei nodi, uso della rete e molto altro. Permette inoltre di generare dei diagrammi inerenti il consumo di risorse (CPU, Memoria usata, tempo d'esecuzione delle build) visualizzandoli in delle \textit{dashboard}. Tali diagrammi sono esportabili in formato ``.\textit{pdf}'';
    \item \textbf{Disk Usage Plugin\footnote{\url{https://plugins.jenkins.io/disk-usage}}}. Crea dei diagrammi e dei report sull'uso della memoria di massa da parte dei singoli \gls{project};
    \item \textbf{Plugin Usage Plugin\footnote{\url{https://wiki.jenkins.io/pages/viewpage.action?pageId=73532011}}}. Mostra in una tabella l'utilizzo dei vari \gls{plugin} di Jenkins, da parte dei vari \gls{project}.
\end{itemize}


%**************************************************************
\section{Backup e Ripristino}
%**************************************************************

Le procedure di backup e ripristino individuate fanno ampio uso del comando \textit{docker cp}\footnote{\url{https://docs.docker.com/engine/reference/commandline/cp/}}. Tale comando permette di copiare file da un \gls{container} verso l'\gls{host} che lo ospita e viceversa. L'uso dei seguenti comandi potrebbe essere integrato in uno script al fine di automatizzare l'intera attività.

\lstinputlisting[language=docker-shell, firstline=18, lastline=20, caption={Sintassi del comando \textit{docker cp} per la creazione di un backup}]{./code/docker-example.txt}  

Nel listing 5.1 è evidenziata la sintassi del comando \textit{docker cp} per copiare la directory principale di Jenkins dal \gls{container} in un determinato percorso del sistema ospite. \\ 

Nel listing 5.2 è invece riportato un esempio completo di backup, salvato in un archivio. Il comando ``\textit{tar}'', con le opzioni \textit{zcvf} permette di:
\begin{itemize}
    \item \textbf{-z}. Comprimere l'archivio usando gzip (solitamente incluso nelle distribuzioni \gls{Linux});
    \item \textbf{-c}. Creare l'archivio;
    \item \textbf{-v}. Mostrare a video i progressi (ed eventuali errori), durante la creazione dell'archivio;
    \item \textbf{-f}. Definire il nome dell'archivio, in questo caso sarà \textit{jenkins-home-AAAAMMDD}, dove
    AAAAMMDD è una nomenclatura rappresentante la data, secondo la seguente convenzione:
    \begin{itemize}
        \item \textbf{AAAA}: Quattro cifre per identificare l'anno;
        \item \textbf{MM}: Due cifre per specificare il mese;
        \item \textbf{GG}: Due cifre per identificare il giorno.
    \end{itemize}
\end{itemize}

\lstinputlisting[language=docker-shell, firstline=22, lastline=28, caption={Esempio completo della creazione di un backup}]{./code/docker-example.txt}  

\lstinputlisting[language=docker-shell, firstline=30, lastline=33, caption={Sintassi del comando \textit{docker cp} per il ripristino}]{./code/docker-example.txt}  

Il procedimento di ripristino seguirà la strada inversa, riportando i dati dal sistema dov'è memorizzato il backup verso il \gls{container}, come evidenziato nel listing 5.3 e tramite un esempio concreto nel listing 5.4 (ipotizzando che la directory principale di Jenkins sia localizzata nel percorso \textit{/opt/datatools/jenkins}).\\ 

\lstinputlisting[language=docker-shell, firstline=35, lastline=37, caption={Sintassi del comando \textit{docker cp} per il ripristino}]{./code/docker-example.txt}  

%**************************************************************
\section{Upgrade}
%**************************************************************

Jenkins ha due linee di sviluppo con distinte tempistiche di rilascio. La prima prevede dei cicli di sviluppo con rilasci settimanali, orientati a fornire nuove funzionalità e \textit{bugfix} rapidi in modo da permettere alla comunità di essere sempre aggiornata con le ultime novità, mentre la seconda denominata \gls{lts} è orientata alla stabilità, con cicli di sviluppo più lunghi e rilascio ogni 12 settimane ed è pensata soprattutto per l'ambito aziendale\footnote{Il ciclo di sviluppo della versione \gls{lts} è riportato in dettaglio nel sito ufficiale di Jenkin \url{https://jenkins.io/download/lts/}}. \\

Il team di sviluppo di Jenkins rilascia il prodotto in vari formati nativi per una variegata tipologia di distribuzioni \gls{Linux} e per Windows; mette inoltre a disposizione delle immagini per \gls{container} già preconfezionate ed infine eroga il prodotto anche sotto forma di applicazione web, in formato ``.\textit{\gls{war}}''. \\

In ambiente containerizzato, la procedura di aggiornamento di Jenkins utilizzando l'archivio ``.\textit{\gls{war}}'' e i volumi di Docker diventa immediata, in dettaglio, basterà associare la directory principale di Jenkins, (da qui in avanti denominata \textit{JENKINS\_HOME}) ad un volume; successivamente basterà sostituire il
file "jenkins.war" precedente, con la nuova versione, nella directory dov'è presente il Dockerfile del
\gls{master} e ricostruire l'\gls{immagine} del \gls{container}.

In seguito, avendo a disposizione un backup, si potrà procedere al ripristino dei dati nel volume. \\
L'aggiornamento dei \gls{plugin}, avviene nello stesso modo.             % Product Prototype
% !TEX encoding = UTF-8
% !TEX TS-program = pdflatex
% !TEX root = ../tesi.tex

%**************************************************************
\pagestyle{IHA-fancy-style}
\chapter{Verifica e validazione}
\label{cap:verifica-validazione}
%**************************************************************

\intro{In questo capitolo vengono riportati i test funzionali eseguiti durante i processi di Verifica e Validazione dell'architettura sviluppata.}

%**************************************************************
\section{Resoconto attività di Testing}
%**************************************************************

In tabella \hyperref[tab:test]{6.1} vengono riportati i test di verifica eseguiti; tali test costituiscono anche una suite di collaudo del sistema. \\

Ciascun test è descritto da:

\begin{itemize}
    \item \textbf{ID}. Un codice alfanumerico che identifica univocamente il test;
    \item \textbf{Descrizione}. Una breve descrizione del test;
    \item \textbf{Comportamento atteso}. Un breve testo che specifica i risultati attesi, in caso di riuscita del test.
    \item \textbf{Stato}. Lo stato d'esecuzione del test, esso può essere:
    \begin{itemize}
        \item \textbf{Riuscito}. Se il test ha avuto esito positivo;
        \item \textbf{Fallito}. Se il test ha avuto esito negativo.
    \end{itemize}
\end{itemize}

\taburowcolors[2] 2{tableLineOne .. tableLineTwo}
\tabulinesep = ^3mm_2mm
\begin{longtabu} to \textwidth {cXXc}
    \caption[Test Funzionali]{Test Funzionali}
    \label{tab:test}
    \endlastfoot
    \rowfont{\bfseries\sffamily\leavevmode\color{white}}
    \rowcolor{tableHeaderBlue}
    %\hline
    \textbf{ID} & \textbf{Descrizione} & \textbf{Comportamento atteso} & \textbf{Stato}\\
    %\hline
    TF001 & Visualizzazione della pagina principale di Jenkins & La pagina principale di Jenkins viene visualizzata & Riuscito\\ %\hline
    TF002 & Autenticazione con l'utente admin di default & L'utente admin riesce ad autenticarsi correttamente con le sue credenziali & Riuscito\\ %\hline
    TF003 & Caricamento dei \gls{project} & Se il container viene rimosso e ne viene creato uno nuovo, con la stessa configurazione i project rimangono & Riuscito\\ %\hline
    TF004 & Caricamento dei \gls{plugin} & Se il container viene rimosso e ne viene creato uno nuovo, con la stessa configurazione i plugin rimangono & Riuscito\\ %\hline
    \newpage
    \rowfont{\bfseries\sffamily\leavevmode\color{white}}
    \rowcolor{tableHeaderBlue}
    %\hline
    \textbf{ID} & \textbf{Descrizione} & \textbf{Comportamento atteso} & \textbf{Stato}\\
    TF005 & \gls{build} di applicazioni web Maven\cite[Chapter~10]{book:maven-reference} & L'applicazione di esempio genera un file ``.\gls{war}'' & Riuscito\\ %\hline
    TF006 & \gls{build} di immagini Docker & L'immagine Docker viene costruita correttamente ed è disponibile nell'\gls{host} & Riuscito\\ %\hline
    TF007 & Esecuzione di \gls{build} distribuite\cite[Chapter~11]{book:jenkins-definitive}\cite{site:jenkins-distributed-builds} con più \gls{host} & Le \gls{build} vengono eseguite correttamente su entrambi i nodi \gls{host} & Riuscito\\ %\hline
    TF008 & Esecuzione di \gls{build} su \gls{slave} \gls{sshg} & Le \gls{build} vengono eseguite correttamente su uno \gls{slave} SSH \gls{container} & Riuscito\\ %\hline
    TF009 & Esecuzione di \gls{build} su \gls{slave} \gls{jnlpg} & Le \gls{build} vengono eseguite correttamente su uno \gls{slave} JNLP \gls{container} & Riuscito\\ %\hline
    TF010 & Recupero degli artefatti dagli \gls{slave} & Gli artefatti vengono copiati correttamente nel \textit{working space} del nodo \gls{master} & Riuscito\\ %\hline
    TF011 & Recupero del codice sorgente da un \gls{repository} \gls{cvs} remoto & La \gls{build} riesce a recuperare il codice sorgente dell'applicazione da costruire da un repository CVS remoto & Riuscito\\ %\hline
    TF012 & Aggiornamento dell'istanza \gls{master} su \gls{container} & L'aggiornamento dell'istanza master avviene correttamente & Riuscito\\ %\hline
    TF013 & Backup dei dati di Jenkins dal volume del \gls{container} \gls{master}& Viene generato un archivio contenente i dati della home directory di Jenkins & Riuscito\\ %\hline
    TF014 & Migrazione dei dati di Jenkins su un altro \gls{container} & La home directory di Jenkins viene trasferita su un nuovo container correttamente & Riuscito\\ %\hline
\end{longtabu}

%**************************************************************
\section{Conclusioni}
%**************************************************************

Tutti i test preventivati hanno riportato esito positivo anche se si evidenziano difficoltà nei test \textit{TF009} e \textit{TF010}. Tali difficoltà sono dovute principalmente a problemi col proxy aziendale.             % Product Design Freeze e SOP
% !TEX encoding = UTF-8
% !TEX TS-program = pdflatex
% !TEX root = ../tesi.tex

%**************************************************************
\pagestyle{IHA-fancy-style}
\chapter{Conclusioni}
\label{cap:conclusioni}
%**************************************************************

\intro{In questo capitolo vengono tratte le considerazioni finali circa le conoscenze acquisite e gli sviluppi futuri del progetto, riportando inoltre una valutazione personale del periodo di stage in rapporto al percorso di studi. Vengono inoltre specificati gli obiettivi raggiunti e i rischi presentatisi.}

%**************************************************************
\section{Raggiungimento degli obiettivi}
\label{sec:obiettivi-raggiunti}
%**************************************************************

Gli obiettivi esposti in \hyperref[sec:ob]{sezione 2.3} e riassunti in \hyperref[tab:obiettivi-iniziali]{tabella 2.1} non anno subito grosse variazioni rispetto quanto inizialmente preventivato. L'obiettivo \textit{O05} relativo allo studio delle modalità di comunicazione tra \gls{master} e \gls{slave}, all'inizio del progetto non era presente e con la sua introduzione si è dovuto ridimensionare l'obiettivo facoltativo \textit{F01} sullo sviluppo d'un \gls{plugin} base. Tale obiettivo facoltativo è stato parzialmente raggiunto anche se non soddisfa pienamente le aspettative aziendali. Non è stato raggiunto l'obiettivo desiderabile \textit{D02} relativo alla stesura di una relazione sui confronti prestazionali tra la piattaforma containerizzata e non, in quanto difficilmente analizzabile in sole otto settimane di stage. Le prestazioni del sistema dipendono da molti parametri (quali ad esempio i limiti imposti ai \gls{container} e le risorse a disposizione dell'\gls{host}). 

Sono stati raggiunti tutti gli obiettivi obbligatori e la quasi totalità di quelli desiderabili. In \hyperref[tab:obiettivi-finali]{tabella 7.1} vengono riassunti gli obiettivi, riportando per ciascuno lo stato di avanzamento, il quale può essere:

\begin{itemize}
    \item \textbf{Raggiunto}. Se l'obiettivo soddisfa pienamente le aspettative aziendali;
    \item \textbf{Parzialmente Raggiunto}. Se l'obiettivo è stato raggiunto ma non è pienamente conforme alle attese;
    \item \textbf{Non Raggiunto}. Se l'obiettivo non ha prodotto i risultati sperati.
\end{itemize}

\newpage
\taburowcolors[2] 2{tableLineOne .. tableLineTwo}
\tabulinesep = ^3mm_2mm
\begin{longtabu} to \textwidth {ccXc}
    \caption[Obiettivi Raggiunti]{Obiettivi Raggiunti}
    \label{tab:obiettivi-finali}
    \endlastfoot
    \rowfont{\bfseries\sffamily\leavevmode\color{white}}
    \rowcolor{darkorange}
    %\hline
    \textbf{ID} & \textbf{Priorità} & \textbf{Descrizione} & Stato \\
    %\hline
    O01 & Obbligatorio & Progettazione dell’architettura di un sistema di \gls{continuous building} & Raggiunto \\ %\hline
    O02 & Obbligatorio & Implementazione dell’architettura progettata su sistemi IKS & Raggiunto \\ %\hline
    O03 & Obbligatorio & Stesura manuale di installazione & Raggiunto \\ %\hline
    O04 & Obbligatorio & Utilizzo della piattaforma implementata per la produzione di \gls{container} Docker & Raggiunto \\ %\hline
    O05 & Obbligatorio & Studio di fattibilità sulle modalità di comunicazione master e slave & Raggiunto \\ %\hline
    D01 & Desiderabile & Containerizzazione della piattaforma implementata & Raggiunto\\ %\hline
    D02 & Desiderabile & Relazione sui confronti prestazionali tra la piattaforma e la sua controparte containerizzata & Non Raggiunto\\ %\hline
    DO3 & Desiderabile & Individuazione metriche per il monitoraggio della piattaforma & Raggiunto \\ %\hline
    DO4 & Desiderabile &  Individuazione procedure per il monitoraggio, backup, ripristino e upgrade dell’infrastruttura & Raggiunto\\ %\hline
    DO5 & Desiderabile & Breve relazione che documenti metriche e procedure dei punti DO3-4 & Raggiunto\\ %\hline
    F01 & Facoltativo & Implementazione di un plugin base per Jenkins & Parzialmente Raggiunto\\ %\hline
\end{longtabu}

Infine, in \hyperref[tab:resume]{tabella 7.2} viene sintetizzata la percentuale di raggiungimento degli obiettivi in base alla priorità associata, sul totale.

\taburowcolors[2] 2{tableLineOne .. tableLineTwo}
\tabulinesep = ^3mm_2mm
\begin{longtabu} to \textwidth {cc}
    \caption[Obiettivi Raggiunti]{Obiettivi Raggiunti}
    \label{tab:resume}
    \endlastfoot
    \rowfont{\bfseries\sffamily\leavevmode\color{white}}
    \rowcolor{darkorange}
    %\hline
    \textbf{Priorità} & \textbf{Percentuale di raggiungimento}\\
    %\hline
    Obbligatori & 100 \% \\ 
    Desiderabili & 80 \% \\ 
    Facoltativi & \textasciitilde100 \% \\ 
\end{longtabu}
\newpage
%**************************************************************
\section{Resoconto dell'analisi dei rischi}
%**************************************************************

In questa sezione si elencano in forma tabellare i rischi precedentemente illustrati in {\hyperref[tab:rischi]{``2.2 - Analisi dei rischi''}} e verificatisi durante lo svolgimento del progetto. Vengono inoltre riportate le misure adottate per la loro risoluzione.

\taburowcolors[2] 2{tableLineOne .. tableLineTwo}
\tabulinesep = ^3mm_2mm
\begin{longtabu} to \textwidth {ccXX}
    \caption[Resoconto dell'analisi dei rischi]{Resoconto dell'analisi dei rischi}
    \label{tab:rischi-final}
    \endlastfoot
    \rowfont{\bfseries\sffamily\leavevmode\color{white}}
    \rowcolor{darkgreen}
    %\hline
    \textbf{Rischio} & \textbf{Verificato} & \textbf{Attualizzazione} & \textbf{Misure adottate} \\
    %\hline
    Scarsa esperienza & NO & L'attività di formazione è stata sufficiente & Nessuna\\ %\hline
    Disponibilità temporali & SI & A causa di esami universitari e piccoli imprevisti alcuni giorni non si è riusciti ad essere in sede. & Il tempo non in sede, non ha intaccato il limite minimo imposto dai vincoli universitari per lo stage curricolare, le attività sono state recuperate lavorando sulla documentazione in remoto. \\ %\hline
    Tecnologie utilizzate & NO & Le tecnologie da utilizzare non si sono rivelate particolarmente complesse & Nessuna\\ %\hline
    Problemi hardware & SI & Si è avuta una temporanea mancanza di connessione, il problema non ha avuto ripercussioni significative & Durante la mancanza di collegamento alla rete esterna ho lavorato offline sulla documentazione \\ %\hline
    Strumenti software & NO & Le componenti non si sono rivelate incompatibili tra loro ma la scarsità di documentazione e la cattiva organizzazione della stessa hanno comportato lievi ritardi. & Oltre alla documentazione ufficiale, ho consultato blog specializzati e testi fisici.\\ %\hline
\end{longtabu}

\newpage

%**************************************************************
\section{Possibili Sviluppi futuri}
%**************************************************************
In questa sezione vengono evidenziati alcuni possibili sviluppi futuri del sistema implementato.\\

\subsection{Monitoraggio}
Come visto in \hyperref[sec:monitoring]{sezione 5.1}, cadvisor non permette nativamente la persistenza dei dati raccolti; inoltre sono presenti altri file di log (per esempio i log delle \gls{build} e i log degli \gls{slave} container) che sono sparsi e di difficile lettura. 

Sarebbe interessante valutare un'integrazione con l'infrastruttura ELK (elasticsearch, Logstash e Kibana) già presente in azienda, soprattutto considerato che cadvisor è integrabile con questo stack tecnologico.

\begin{figure}[H]
    \capstart
    \centering
    \includegraphics[width=7cm, keepaspectratio]{../immagini/elk.png}
    \caption{Loghi dello stack ELK}
\end{figure}

\subsection{Repository prodotti}

Scopo dello stage era l'implementazione di tutto il sistema di \gls{continuous building} ma non è mai stato fatto un vero uso degli artefatti prodotti. Un possibile sviluppo sarebbe integrare la soluzione da me sviluppata con un \gls{repository} prodotti come Nexus Repository Manager\footnote{\url{https://www.sonatype.com/nexus-repository-sonatype}}, già utilizzato con profitto dall'azienda.\\

Anche per le immagini Docker sarebbe utile approfondire l'integrazione del sistema con un \textit{registry} pubblico o privato, cioè con un repository appositamente pensato per la memorizzazione, la gestione e la distribuzione di immagini Docker.

%**************************************************************
\section{Conoscenze acquisite}
%**************************************************************

L'attività di stage mi ha permesso di apprendere nuove tecnologie e rapportarmi con aspetti diversi da quelli affrontati durante gli anni di studio, soprattutto ho avuto modo di seguire un progetto dalle sue fasi iniziali a quelle finali in una realtà aziendale. Nello specifico, durante il periodo di stage:

\begin{itemize}
    \item Ho approfondito le mie conoscenze del sistema operativo \gls{Linux}, della sua struttura e delle sue logiche di funzionamento; soprattutto nella gestione di utenti, gruppi e relativi permessi. Ho compreso come configurare servizi tramite \gls{systemd} in modo da eseguirli, fermarli, abilitarli in automatico con l'avvio del sistema ed infine ho sviluppato conoscenze sulle architetture distribuite e containerizzate in una rete\cite{site:docker-networking}; 
    \item Ho ottenuto competenze nelle tecnologie di containerizzazione, in particolare nel creare e gestire immagini Docker, nel configurare container e nel redigere i Dockerfile;
    \item Ho avuto modo di rapportarmi con metodologie di sviluppo mai affrontate e da me poco conosciute;
    \item Ho approfondito gli aspetti dell'\textit{integrazione continua} e dell'automazione di compiti ripetibili.
    \item Ho migliorato le mie capacità di \textit{problem solving} adattando tecnologie open source ai bisogni dell'azienda.
\end{itemize}

%**************************************************************
\section{Valutazione personale}
%**************************************************************

Nel complesso, seppur con le molte difficoltà affrontate a livello logistico e le tempistiche molto ristrette valuto positivamente l'esperienza di stage. \\
Le mie motivazioni personali sono state pienamente soddisfatte e anche se all'inizio avevo conoscenze quasi nulle del dominio applicativo e delle tecnologie da utilizzare, gli insegnamenti forniti e l'ampio bagaglio teorico e metodologico fornito dal corso di studi mi hanno permesso di riuscire comunque a raggiungere gli obiettivi prefissati.\\
Dal punto di vista professionale, non è ben chiaro se vi sia possibilità di proseguire con il progetto o in questo ambito lavorativo all'interno dell'azienda ospitante, comunque il bagaglio di conoscenze acquisite e l'esperienza maturata saranno sicuramente utili anche in altri contesti lavorativi.             % Conclusioni
\appendix                               
% !TEX encoding = UTF-8
% !TEX TS-program = pdflatex
% !TEX root = ../tesi.tex

%**************************************************************
\chapter{Appendice A}
%**************************************************************

\epigraph{Citazione}{Autore della citazione}



             % Appendice A

%**************************************************************
% Materiale finale
%**************************************************************
%\backmatter
%\printglossaries
%% !TEX encoding = UTF-8
% !TEX TS-program = pdflatex
% !TEX root = ../tesi.tex

%**************************************************************
% Bibliografia
%**************************************************************

\cleardoublepage
\chapter{Bibliografia}

\nocite{*}
% Stampa i riferimenti bibliografici
\printbibliography[heading=subbibliography,title={Riferimenti bibliografici},type=book]

% Stampa i siti web consultati
\printbibliography[heading=subbibliography,title={Siti web consultati},type=online]


\end{document}